\documentclass[11pt]{report}

% Shared preamble for main.tex and Proof_Attempts.tex
% (packages, theorem styles, macros, page layout, etc.)

% -------------------------------------------------
% Page layout / typography
% -------------------------------------------------
\usepackage[a4paper,margin=30mm,headheight=14pt]{geometry}
\usepackage{microtype}              % better spacing/kerning
\usepackage[T1]{fontenc}
\usepackage[utf8]{inputenc}         % ok for pdflatex; remove if using lualatex/xelatex
\usepackage{xcolor}
\usepackage{lmodern}                % modern Latin Modern fonts
\usepackage{setspace}               % optional: \onehalfspacing etc.

% -------------------------------------------------
% Graphics / tables / lists
% -------------------------------------------------
\usepackage{graphicx}
\usepackage{booktabs}               % professional tables
\usepackage{array}
\usepackage{multirow}
\usepackage{caption}
\usepackage{subcaption}
\usepackage{enumitem}

% -------------------------------------------------
% Math (standard praxis)
% -------------------------------------------------
\usepackage{amsmath,amssymb,amsthm} % AMS core
\usepackage{mathtools}             % extends/fixes amsmath
\usepackage{stmaryrd}              % \llbracket, \rrbracket for stochastic intervals

% Upright differential (use as \int f(x)\,\diff x)
\newcommand{\diff}{\mathop{}\!\mathrm{d}}

% Operators
\DeclareMathOperator*{\arginf}{arginf}
\DeclareMathOperator*{\argmin}{arg\,min}

\usepackage{bm}                    % bold math symbols
\usepackage{dsfont}                % \mathds{1} indicator
\usepackage{mathrsfs}              % script font \mathscr
\usepackage{bbm}                   % alternative blackboard bold, incl. \mathbbm{1}
\usepackage{siunitx}               % units/numbers (also nice in finance tables)

% -------------------------------------------------
% References / hyperlinks (standard modern setup)
% -------------------------------------------------
\usepackage{xcolor}
\usepackage{hyperref}
\hypersetup{
  colorlinks=true,
  linkcolor=blue,
  citecolor=blue,
  urlcolor=blue,
  pdftitle={Master Project},
  pdfauthor={Your Name},
  pdfsubject={Mathematical Finance},
  pdfkeywords={Mathematical Finance, Stochastic Calculus, ...}
}

% cleveref + shared theorem counters:
% If environments share the "theorem" counter, \cref may otherwise call everything "Theorem".
% aliascnt creates alias counters so cleveref can distinguish Remark/Lemma/etc.
\usepackage{aliascnt}

% cleveref must be loaded AFTER hyperref
\usepackage[nameinlink,noabbrev]{cleveref}

% -------------------------------------------------
% Theorems / definitions (standard mathematical setup)
% -------------------------------------------------
% Numbering: Theorem 2.3, Lemma 2.4, ... within each chapter.
\theoremstyle{plain} % italic body, bold heading
\newtheorem{theorem}{Theorem}[chapter]

\newaliascnt{proposition}{theorem}
\newtheorem{proposition}[proposition]{Proposition}
\aliascntresetthe{proposition}

\newaliascnt{lemma}{theorem}
\newtheorem{lemma}[lemma]{Lemma}
\aliascntresetthe{lemma}

\newaliascnt{corollary}{theorem}
\newtheorem{corollary}[corollary]{Corollary}
\aliascntresetthe{corollary}

\theoremstyle{definition} % upright body
\newaliascnt{definition}{theorem}
\newtheorem{definition}[definition]{Definition}
\aliascntresetthe{definition}

\newaliascnt{assumption}{theorem}
\newtheorem{assumption}[assumption]{Assumption}
\aliascntresetthe{assumption}

\theoremstyle{remark} % upright body, italic heading
\newaliascnt{remark}{theorem}
\newtheorem{remark}[remark]{Remark}
\aliascntresetthe{remark}

% Unnumbered variants (use sparingly)
\newtheorem*{theorem*}{Theorem}
\newtheorem*{lemma*}{Lemma}
\newtheorem*{definition*}{Definition}

% Cleveref names for theorem-like environments
\crefname{theorem}{Theorem}{Theorems}
\Crefname{theorem}{Theorem}{Theorems}
\crefname{proposition}{Proposition}{Propositions}
\Crefname{proposition}{Proposition}{Propositions}
\crefname{lemma}{Lemma}{Lemmas}
\Crefname{lemma}{Lemma}{Lemmas}
\crefname{corollary}{Corollary}{Corollaries}
\Crefname{corollary}{Corollary}{Corollaries}
\crefname{definition}{Definition}{Definitions}
\Crefname{definition}{Definition}{Definitions}
\crefname{assumption}{Assumption}{Assumptions}
\Crefname{assumption}{Assumption}{Assumptions}
\crefname{remark}{Remark}{Remarks}
\Crefname{remark}{Remark}{Remarks}

% -------------------------------------------------
% Bibliography (biblatex + biber)
% -------------------------------------------------
\usepackage{csquotes} % recommended with biblatex
\usepackage[
  backend=biber,
  style=authoryear,
  dashed=false, % don't replace repeated author/editor with an em-dash
  maxcitenames=2,
  maxbibnames=99
]{biblatex}
\addbibresource{references.bib}

% -------------------------------------------------
% Appendix + PDF inclusion (optional)
% -------------------------------------------------
\usepackage{appendix}
\usepackage{pdfpages}

% -------------------------------------------------
% Nomenclature / symbols list (optional but common in math finance)
% -------------------------------------------------
\usepackage[intoc]{nomencl}
\makenomenclature
% Example: \nomenclature{$W_t$}{Standard Brownian motion}

% -------------------------------------------------
% Header / footer (simple, thesis-friendly)
% -------------------------------------------------
\usepackage{fancyhdr}
\pagestyle{fancy}
\fancyhf{}
\fancyhead[LE,RO]{\thepage}
\fancyhead[LO]{\nouppercase{\rightmark}}
\fancyhead[RE]{\nouppercase{\leftmark}}
\renewcommand{\headrulewidth}{0.4pt}

% -------------------------------------------------
% Table of contents depth (DETAILED ToC)
% -------------------------------------------------
% tocdepth: 0=part, 1=chapter, 2=section, 3=subsection, 4=subsubsection
\setcounter{tocdepth}{3}  % show down to \subsection in ToC
\setcounter{secnumdepth}{3} % number down to \subsection (or 4 for subsubsection)

\begin{document}

% This file is included from main.tex via \documentclass[11pt]{report}

% Shared preamble for main.tex and Proof_Attempts.tex
% (packages, theorem styles, macros, page layout, etc.)

% -------------------------------------------------
% Page layout / typography
% -------------------------------------------------
\usepackage[a4paper,margin=30mm,headheight=14pt]{geometry}
\usepackage{microtype}              % better spacing/kerning
\usepackage[T1]{fontenc}
\usepackage[utf8]{inputenc}         % ok for pdflatex; remove if using lualatex/xelatex
\usepackage{xcolor}
\usepackage{lmodern}                % modern Latin Modern fonts
\usepackage{setspace}               % optional: \onehalfspacing etc.

% -------------------------------------------------
% Graphics / tables / lists
% -------------------------------------------------
\usepackage{graphicx}
\usepackage{booktabs}               % professional tables
\usepackage{array}
\usepackage{multirow}
\usepackage{caption}
\usepackage{subcaption}
\usepackage{enumitem}

% -------------------------------------------------
% Math (standard praxis)
% -------------------------------------------------
\usepackage{amsmath,amssymb,amsthm} % AMS core
\usepackage{mathtools}             % extends/fixes amsmath
\usepackage{stmaryrd}              % \llbracket, \rrbracket for stochastic intervals

% Upright differential (use as \int f(x)\,\diff x)
\newcommand{\diff}{\mathop{}\!\mathrm{d}}

% Operators
\DeclareMathOperator*{\arginf}{arginf}
\DeclareMathOperator*{\argmin}{arg\,min}

\usepackage{bm}                    % bold math symbols
\usepackage{dsfont}                % \mathds{1} indicator
\usepackage{mathrsfs}              % script font \mathscr
\usepackage{bbm}                   % alternative blackboard bold, incl. \mathbbm{1}
\usepackage{siunitx}               % units/numbers (also nice in finance tables)

% -------------------------------------------------
% References / hyperlinks (standard modern setup)
% -------------------------------------------------
\usepackage{xcolor}
\usepackage{hyperref}
\hypersetup{
  colorlinks=true,
  linkcolor=blue,
  citecolor=blue,
  urlcolor=blue,
  pdftitle={Master Project},
  pdfauthor={Your Name},
  pdfsubject={Mathematical Finance},
  pdfkeywords={Mathematical Finance, Stochastic Calculus, ...}
}

% cleveref + shared theorem counters:
% If environments share the "theorem" counter, \cref may otherwise call everything "Theorem".
% aliascnt creates alias counters so cleveref can distinguish Remark/Lemma/etc.
\usepackage{aliascnt}

% cleveref must be loaded AFTER hyperref
\usepackage[nameinlink,noabbrev]{cleveref}

% -------------------------------------------------
% Theorems / definitions (standard mathematical setup)
% -------------------------------------------------
% Numbering: Theorem 2.3, Lemma 2.4, ... within each chapter.
\theoremstyle{plain} % italic body, bold heading
\newtheorem{theorem}{Theorem}[chapter]

\newaliascnt{proposition}{theorem}
\newtheorem{proposition}[proposition]{Proposition}
\aliascntresetthe{proposition}

\newaliascnt{lemma}{theorem}
\newtheorem{lemma}[lemma]{Lemma}
\aliascntresetthe{lemma}

\newaliascnt{corollary}{theorem}
\newtheorem{corollary}[corollary]{Corollary}
\aliascntresetthe{corollary}

\theoremstyle{definition} % upright body
\newaliascnt{definition}{theorem}
\newtheorem{definition}[definition]{Definition}
\aliascntresetthe{definition}

\newaliascnt{assumption}{theorem}
\newtheorem{assumption}[assumption]{Assumption}
\aliascntresetthe{assumption}

\theoremstyle{remark} % upright body, italic heading
\newaliascnt{remark}{theorem}
\newtheorem{remark}[remark]{Remark}
\aliascntresetthe{remark}

% Unnumbered variants (use sparingly)
\newtheorem*{theorem*}{Theorem}
\newtheorem*{lemma*}{Lemma}
\newtheorem*{definition*}{Definition}

% Cleveref names for theorem-like environments
\crefname{theorem}{Theorem}{Theorems}
\Crefname{theorem}{Theorem}{Theorems}
\crefname{proposition}{Proposition}{Propositions}
\Crefname{proposition}{Proposition}{Propositions}
\crefname{lemma}{Lemma}{Lemmas}
\Crefname{lemma}{Lemma}{Lemmas}
\crefname{corollary}{Corollary}{Corollaries}
\Crefname{corollary}{Corollary}{Corollaries}
\crefname{definition}{Definition}{Definitions}
\Crefname{definition}{Definition}{Definitions}
\crefname{assumption}{Assumption}{Assumptions}
\Crefname{assumption}{Assumption}{Assumptions}
\crefname{remark}{Remark}{Remarks}
\Crefname{remark}{Remark}{Remarks}

% -------------------------------------------------
% Bibliography (biblatex + biber)
% -------------------------------------------------
\usepackage{csquotes} % recommended with biblatex
\usepackage[
  backend=biber,
  style=authoryear,
  dashed=false, % don't replace repeated author/editor with an em-dash
  maxcitenames=2,
  maxbibnames=99
]{biblatex}
\addbibresource{references.bib}

% -------------------------------------------------
% Appendix + PDF inclusion (optional)
% -------------------------------------------------
\usepackage{appendix}
\usepackage{pdfpages}

% -------------------------------------------------
% Nomenclature / symbols list (optional but common in math finance)
% -------------------------------------------------
\usepackage[intoc]{nomencl}
\makenomenclature
% Example: \nomenclature{$W_t$}{Standard Brownian motion}

% -------------------------------------------------
% Header / footer (simple, thesis-friendly)
% -------------------------------------------------
\usepackage{fancyhdr}
\pagestyle{fancy}
\fancyhf{}
\fancyhead[LE,RO]{\thepage}
\fancyhead[LO]{\nouppercase{\rightmark}}
\fancyhead[RE]{\nouppercase{\leftmark}}
\renewcommand{\headrulewidth}{0.4pt}

% -------------------------------------------------
% Table of contents depth (DETAILED ToC)
% -------------------------------------------------
% tocdepth: 0=part, 1=chapter, 2=section, 3=subsection, 4=subsubsection
\setcounter{tocdepth}{3}  % show down to \subsection in ToC
\setcounter{secnumdepth}{3} % number down to \subsection (or 4 for subsubsection)

\begin{document}

% This file is included from main.tex via \documentclass[11pt]{report}

% Shared preamble for main.tex and Proof_Attempts.tex
% (packages, theorem styles, macros, page layout, etc.)

% -------------------------------------------------
% Page layout / typography
% -------------------------------------------------
\usepackage[a4paper,margin=30mm,headheight=14pt]{geometry}
\usepackage{microtype}              % better spacing/kerning
\usepackage[T1]{fontenc}
\usepackage[utf8]{inputenc}         % ok for pdflatex; remove if using lualatex/xelatex
\usepackage{xcolor}
\usepackage{lmodern}                % modern Latin Modern fonts
\usepackage{setspace}               % optional: \onehalfspacing etc.

% -------------------------------------------------
% Graphics / tables / lists
% -------------------------------------------------
\usepackage{graphicx}
\usepackage{booktabs}               % professional tables
\usepackage{array}
\usepackage{multirow}
\usepackage{caption}
\usepackage{subcaption}
\usepackage{enumitem}

% -------------------------------------------------
% Math (standard praxis)
% -------------------------------------------------
\usepackage{amsmath,amssymb,amsthm} % AMS core
\usepackage{mathtools}             % extends/fixes amsmath
\usepackage{stmaryrd}              % \llbracket, \rrbracket for stochastic intervals

% Upright differential (use as \int f(x)\,\diff x)
\newcommand{\diff}{\mathop{}\!\mathrm{d}}

% Operators
\DeclareMathOperator*{\arginf}{arginf}
\DeclareMathOperator*{\argmin}{arg\,min}

\usepackage{bm}                    % bold math symbols
\usepackage{dsfont}                % \mathds{1} indicator
\usepackage{mathrsfs}              % script font \mathscr
\usepackage{bbm}                   % alternative blackboard bold, incl. \mathbbm{1}
\usepackage{siunitx}               % units/numbers (also nice in finance tables)

% -------------------------------------------------
% References / hyperlinks (standard modern setup)
% -------------------------------------------------
\usepackage{xcolor}
\usepackage{hyperref}
\hypersetup{
  colorlinks=true,
  linkcolor=blue,
  citecolor=blue,
  urlcolor=blue,
  pdftitle={Master Project},
  pdfauthor={Your Name},
  pdfsubject={Mathematical Finance},
  pdfkeywords={Mathematical Finance, Stochastic Calculus, ...}
}

% cleveref + shared theorem counters:
% If environments share the "theorem" counter, \cref may otherwise call everything "Theorem".
% aliascnt creates alias counters so cleveref can distinguish Remark/Lemma/etc.
\usepackage{aliascnt}

% cleveref must be loaded AFTER hyperref
\usepackage[nameinlink,noabbrev]{cleveref}

% -------------------------------------------------
% Theorems / definitions (standard mathematical setup)
% -------------------------------------------------
% Numbering: Theorem 2.3, Lemma 2.4, ... within each chapter.
\theoremstyle{plain} % italic body, bold heading
\newtheorem{theorem}{Theorem}[chapter]

\newaliascnt{proposition}{theorem}
\newtheorem{proposition}[proposition]{Proposition}
\aliascntresetthe{proposition}

\newaliascnt{lemma}{theorem}
\newtheorem{lemma}[lemma]{Lemma}
\aliascntresetthe{lemma}

\newaliascnt{corollary}{theorem}
\newtheorem{corollary}[corollary]{Corollary}
\aliascntresetthe{corollary}

\theoremstyle{definition} % upright body
\newaliascnt{definition}{theorem}
\newtheorem{definition}[definition]{Definition}
\aliascntresetthe{definition}

\newaliascnt{assumption}{theorem}
\newtheorem{assumption}[assumption]{Assumption}
\aliascntresetthe{assumption}

\theoremstyle{remark} % upright body, italic heading
\newaliascnt{remark}{theorem}
\newtheorem{remark}[remark]{Remark}
\aliascntresetthe{remark}

% Unnumbered variants (use sparingly)
\newtheorem*{theorem*}{Theorem}
\newtheorem*{lemma*}{Lemma}
\newtheorem*{definition*}{Definition}

% Cleveref names for theorem-like environments
\crefname{theorem}{Theorem}{Theorems}
\Crefname{theorem}{Theorem}{Theorems}
\crefname{proposition}{Proposition}{Propositions}
\Crefname{proposition}{Proposition}{Propositions}
\crefname{lemma}{Lemma}{Lemmas}
\Crefname{lemma}{Lemma}{Lemmas}
\crefname{corollary}{Corollary}{Corollaries}
\Crefname{corollary}{Corollary}{Corollaries}
\crefname{definition}{Definition}{Definitions}
\Crefname{definition}{Definition}{Definitions}
\crefname{assumption}{Assumption}{Assumptions}
\Crefname{assumption}{Assumption}{Assumptions}
\crefname{remark}{Remark}{Remarks}
\Crefname{remark}{Remark}{Remarks}

% -------------------------------------------------
% Bibliography (biblatex + biber)
% -------------------------------------------------
\usepackage{csquotes} % recommended with biblatex
\usepackage[
  backend=biber,
  style=authoryear,
  dashed=false, % don't replace repeated author/editor with an em-dash
  maxcitenames=2,
  maxbibnames=99
]{biblatex}
\addbibresource{references.bib}

% -------------------------------------------------
% Appendix + PDF inclusion (optional)
% -------------------------------------------------
\usepackage{appendix}
\usepackage{pdfpages}

% -------------------------------------------------
% Nomenclature / symbols list (optional but common in math finance)
% -------------------------------------------------
\usepackage[intoc]{nomencl}
\makenomenclature
% Example: \nomenclature{$W_t$}{Standard Brownian motion}

% -------------------------------------------------
% Header / footer (simple, thesis-friendly)
% -------------------------------------------------
\usepackage{fancyhdr}
\pagestyle{fancy}
\fancyhf{}
\fancyhead[LE,RO]{\thepage}
\fancyhead[LO]{\nouppercase{\rightmark}}
\fancyhead[RE]{\nouppercase{\leftmark}}
\renewcommand{\headrulewidth}{0.4pt}

% -------------------------------------------------
% Table of contents depth (DETAILED ToC)
% -------------------------------------------------
% tocdepth: 0=part, 1=chapter, 2=section, 3=subsection, 4=subsubsection
\setcounter{tocdepth}{3}  % show down to \subsection in ToC
\setcounter{secnumdepth}{3} % number down to \subsection (or 4 for subsubsection)

\begin{document}

% This file is included from main.tex via \documentclass[11pt]{report}

% Shared preamble for main.tex and Proof_Attempts.tex
% (packages, theorem styles, macros, page layout, etc.)

% -------------------------------------------------
% Page layout / typography
% -------------------------------------------------
\usepackage[a4paper,margin=30mm,headheight=14pt]{geometry}
\usepackage{microtype}              % better spacing/kerning
\usepackage[T1]{fontenc}
\usepackage[utf8]{inputenc}         % ok for pdflatex; remove if using lualatex/xelatex
\usepackage{xcolor}
\usepackage{lmodern}                % modern Latin Modern fonts
\usepackage{setspace}               % optional: \onehalfspacing etc.

% -------------------------------------------------
% Graphics / tables / lists
% -------------------------------------------------
\usepackage{graphicx}
\usepackage{booktabs}               % professional tables
\usepackage{array}
\usepackage{multirow}
\usepackage{caption}
\usepackage{subcaption}
\usepackage{enumitem}

% -------------------------------------------------
% Math (standard praxis)
% -------------------------------------------------
\usepackage{amsmath,amssymb,amsthm} % AMS core
\usepackage{mathtools}             % extends/fixes amsmath
\usepackage{stmaryrd}              % \llbracket, \rrbracket for stochastic intervals

% Upright differential (use as \int f(x)\,\diff x)
\newcommand{\diff}{\mathop{}\!\mathrm{d}}

% Operators
\DeclareMathOperator*{\arginf}{arginf}
\DeclareMathOperator*{\argmin}{arg\,min}

\usepackage{bm}                    % bold math symbols
\usepackage{dsfont}                % \mathds{1} indicator
\usepackage{mathrsfs}              % script font \mathscr
\usepackage{bbm}                   % alternative blackboard bold, incl. \mathbbm{1}
\usepackage{siunitx}               % units/numbers (also nice in finance tables)

% -------------------------------------------------
% References / hyperlinks (standard modern setup)
% -------------------------------------------------
\usepackage{xcolor}
\usepackage{hyperref}
\hypersetup{
  colorlinks=true,
  linkcolor=blue,
  citecolor=blue,
  urlcolor=blue,
  pdftitle={Master Project},
  pdfauthor={Your Name},
  pdfsubject={Mathematical Finance},
  pdfkeywords={Mathematical Finance, Stochastic Calculus, ...}
}

% cleveref + shared theorem counters:
% If environments share the "theorem" counter, \cref may otherwise call everything "Theorem".
% aliascnt creates alias counters so cleveref can distinguish Remark/Lemma/etc.
\usepackage{aliascnt}

% cleveref must be loaded AFTER hyperref
\usepackage[nameinlink,noabbrev]{cleveref}

% -------------------------------------------------
% Theorems / definitions (standard mathematical setup)
% -------------------------------------------------
% Numbering: Theorem 2.3, Lemma 2.4, ... within each chapter.
\theoremstyle{plain} % italic body, bold heading
\newtheorem{theorem}{Theorem}[chapter]

\newaliascnt{proposition}{theorem}
\newtheorem{proposition}[proposition]{Proposition}
\aliascntresetthe{proposition}

\newaliascnt{lemma}{theorem}
\newtheorem{lemma}[lemma]{Lemma}
\aliascntresetthe{lemma}

\newaliascnt{corollary}{theorem}
\newtheorem{corollary}[corollary]{Corollary}
\aliascntresetthe{corollary}

\theoremstyle{definition} % upright body
\newaliascnt{definition}{theorem}
\newtheorem{definition}[definition]{Definition}
\aliascntresetthe{definition}

\newaliascnt{assumption}{theorem}
\newtheorem{assumption}[assumption]{Assumption}
\aliascntresetthe{assumption}

\theoremstyle{remark} % upright body, italic heading
\newaliascnt{remark}{theorem}
\newtheorem{remark}[remark]{Remark}
\aliascntresetthe{remark}

% Unnumbered variants (use sparingly)
\newtheorem*{theorem*}{Theorem}
\newtheorem*{lemma*}{Lemma}
\newtheorem*{definition*}{Definition}

% Cleveref names for theorem-like environments
\crefname{theorem}{Theorem}{Theorems}
\Crefname{theorem}{Theorem}{Theorems}
\crefname{proposition}{Proposition}{Propositions}
\Crefname{proposition}{Proposition}{Propositions}
\crefname{lemma}{Lemma}{Lemmas}
\Crefname{lemma}{Lemma}{Lemmas}
\crefname{corollary}{Corollary}{Corollaries}
\Crefname{corollary}{Corollary}{Corollaries}
\crefname{definition}{Definition}{Definitions}
\Crefname{definition}{Definition}{Definitions}
\crefname{assumption}{Assumption}{Assumptions}
\Crefname{assumption}{Assumption}{Assumptions}
\crefname{remark}{Remark}{Remarks}
\Crefname{remark}{Remark}{Remarks}

% -------------------------------------------------
% Bibliography (biblatex + biber)
% -------------------------------------------------
\usepackage{csquotes} % recommended with biblatex
\usepackage[
  backend=biber,
  style=authoryear,
  dashed=false, % don't replace repeated author/editor with an em-dash
  maxcitenames=2,
  maxbibnames=99
]{biblatex}
\addbibresource{references.bib}

% -------------------------------------------------
% Appendix + PDF inclusion (optional)
% -------------------------------------------------
\usepackage{appendix}
\usepackage{pdfpages}

% -------------------------------------------------
% Nomenclature / symbols list (optional but common in math finance)
% -------------------------------------------------
\usepackage[intoc]{nomencl}
\makenomenclature
% Example: \nomenclature{$W_t$}{Standard Brownian motion}

% -------------------------------------------------
% Header / footer (simple, thesis-friendly)
% -------------------------------------------------
\usepackage{fancyhdr}
\pagestyle{fancy}
\fancyhf{}
\fancyhead[LE,RO]{\thepage}
\fancyhead[LO]{\nouppercase{\rightmark}}
\fancyhead[RE]{\nouppercase{\leftmark}}
\renewcommand{\headrulewidth}{0.4pt}

% -------------------------------------------------
% Table of contents depth (DETAILED ToC)
% -------------------------------------------------
% tocdepth: 0=part, 1=chapter, 2=section, 3=subsection, 4=subsubsection
\setcounter{tocdepth}{3}  % show down to \subsection in ToC
\setcounter{secnumdepth}{3} % number down to \subsection (or 4 for subsubsection)

\begin{document}

% This file is included from main.tex via \input{Proof_Attempts}
% so it automatically matches the exact same layout/preamble.

\section{Proof attempts}

% Put your drafts here. Example skeleton:
%
%\begin{proof}[Proof attempt (Lemma~\ref{lem:...})]
%    ...
%\end{proof}
\section{Extension to semi-martingale case}
\cref{thm:cont-deep-ms-hedging-martingale} was in the realm of classical martingale stochastic calculus, with integrators \(M\in\mathcal{H}^{2}\)(even \(\mathcal{H}^{2+\epsilon}\) and integrands \(\vartheta\in L^2(M)\). In light of \cref{thm:Approx of stoch. Integrals}, we aim to prove a similar result as in \cref{thm:cont-deep-ms-hedging-martingale} for the more general case where \(M\) is a semimartingale. Following \parencite{Schweizer2001GuidedTourQuadraticHedging} we restrict the class of possible strategies to \(\vartheta_{GLP}\) such that \(G_T(\vartheta_{GLP})\) is a closed subspace in \(L^2(\mathbb{P})\). In this case, we know that the minimization problem has a solution, as \(\mathcal{A}=\mathbb{R}+G_T(\Theta_{GLP})\) is a closed subspace in \(L^2(\mathbb{P})\).; this means there exists an optimal trading strategy \((V_0^*,\vartheta^*)\in\mathcal{A}=\mathbb{R}+G_T(\Theta_{GLP})\) (\textcolor{blue}{I guess we don't even need strong conditions to get \(\mathcal{A}\) closed as we don't need explicit expression of the optimal trading strategy at first, hence we can directly project on \(\bar{\mathcal{A}}\) which is closed in \(L^2(\mathbb{P})\). In that line of thought, any \(\Theta\subset\Theta_2\) is sufficient for the existence of a solution in \(\bar{\mathcal{A}}\)}.) Then we have the following consequence of the Hilbert Projection theorem
\begin{theorem}[Hilbert Projection for Mean Sqaure Hedging]\label{thm:Hilbert Projection for Mean Sqaure Hedging}
    Let \(\Theta\subset\Theta_2\) and \(\mathcal{A}=\mathbb{R}+G_T(\Theta)\subset L^2(\mathbb{P})\). Then the \(L^2\)-closure \(\bar{\mathcal{A}}\) is a closed subspace of \(L^2(\mathbb{P})\). Hence, there exists \((V_0^*,\vartheta^*)\) such that 
    \begin{equation}
   \lVert H-V_0^*-\int_{0}^T\vartheta^*_s\diff X_s\rVert_{L^2}= \inf_{(c,\vartheta)\in\bar{\mathcal{A}}}\lVert H-c-\int_{0}^T\vartheta_s\diff X_s\rVert_{L^2}
    \end{equation}
    and the solution is unique. 
\end{theorem}
\begin{remark}
    Notice, we can study the related problem where we fix \(V_0\) (corresponds to the situation where the price \(V_0\) of the liability \(H\) might be determined by the market) and study
    \begin{equation}
        \inf_{\vartheta\in\Theta}\lVert H-V_0-\int_{0}^T\vartheta_s\diff X_s\rVert_{L^2}
    \end{equation}
    Again we can consider the the \(L^2\) closure of \(G_T(\Theta)=\left\{\int_0^T \vartheta_s\diff X_s:\vartheta\in\Theta\right\}\) and use the Hilbert Projection theorem on this closed subspace of \(L^2\).
\end{remark}
Then, crucially, notice that from \cref{thm:Approx of stoch. Integrals} we obtain the convergence in probability of the final values of the respective integral processes. Indeed, the \(\mathbb{P}\)-\(\mathbb{P}\)-ucp convergence of the integral processes, restricted to a finite time horizon \(T\), reads as:
\begin{equation}
    \sup_{0\leq t\leq T}\left|\int_{0}^t\vartheta_s^n\diff X_s-\int_{0}^t\vartheta^*_s\diff X_s\right|\to 0\quad\text{in Probability}
\end{equation}
where \(\{\vartheta^n\}_{n\in\mathbb{N}}\in\mathcal{S}(\psi)\) is a minimizing sequence from \cref{thm:Approx of stoch. Integrals}
\\\\
In particular, this implies 
\begin{equation}
    \int_{0}^T\vartheta_s^n\diff X_s\xrightarrow{\mathbb{P}}\int_{0}^T\vartheta^*_s\diff X_s\;,
\end{equation}
where convergence holds in probability. Imposing the additional assumption that \(\left\{\left|\int_{0}^T\vartheta_s^n\diff X_s\right|^2\right\}_{n\in\mathbb{N}}\) is a uniformly integrable family of random variables  leads to the \(L^2\) convergence\footnote{In particular, the condition also implies that the random variable \(\int_{0}^T\vartheta_s^n\diff X_s\in L^2(\mathbb{P)}\quad\forall n\in\mathbb{N}\)}: 
\begin{equation}
\int_{0}^T\vartheta_s^n\diff X_s \xrightarrow{L^2(\mathbb{P})} \int_{0}^T\vartheta^*_s\diff X_s\;.
\end{equation}
We do NOT want the uniform integrability of \(\left\{\left|\int_{0}^T\vartheta_s^n\diff X_s\right|^2\right\}_{n\in\mathbb{N}}\) to follow from conditions on a sequence of trading strategies. Indeed, \cref{thm:Approx of stoch. Integrals} asserts the existence of a sequence \(\{\vartheta^n\}_{n\in\mathbb{N}}\subset\mathcal{S}(\psi)\). We don't want a theorem that says something like: 
\\\\
"If that sequence happens to be square uniformly integrable, then it follows...."
\\\\
Hence, if necessary, we shall impose conditions only on the integrator \(X\). Recall that for \(\vartheta_n\in\mathcal{S}(\psi)\) that is \(\vartheta^n_t=\varphi_0\mathbbm{1}_{\{0\}}+\sum_{i=1}^l\varphi_i^n\mathbbm{1}_{\tau_i,\tau_{i+1}}\) and \(\varphi_i^{(n)}=f^{(n)}(Y^{\tau_i}_{t_1},\dots,Y^{\tau_i}_{t_r})\) where \(f\in\mathcal{NN}((\mathbb{R}^{d\times r})^*,\psi)\)\footnote{the parameter \(r\in\mathbb{N}\) is not further specified. It relies to the existence result of \parencite[Thm.~5.23.]{Arandjelovic2025}}, the stochastic integral takes the form \(\int_0^T\vartheta^n_s\diff X_s = \sum_{i=1}^k \varphi^{(n)}_i(X_{\tau_i\wedge T}-X_{\tau_{i+1}\wedge T})\). Note that \(\left|\int_0^T\vartheta^n_s\diff X_s\right|^2 \leq \sum_{i=1}^k \left|\varphi^{(n)}_i\right|^2\left|(X_{\tau_i\wedge T}-X_{\tau_{i+1}\wedge T})\right|^2\). But there can't exist a uniform bound (in n) on \(|\varphi_i^n|\) as this would imply that neural networks can't approximate constants above this bound which is absurd. Also, notice for \(\tau<\sigma\) two stopping times, that \(f_n = c_n\mathbbm{1}_{[\tau,\sigma)}\in\mathcal{S}(\psi)\) for all \(c_n\in\mathbb{R}\). Take \(c_n\to\infty\) as \(n\to\infty\) a real deterministic sequence, then \(\int_0^T(f_n)_s\diff X_s=c_n(X_{\tau\wedge T}-X_{\sigma\wedge T})\), and hence \(\mathbb{E}[\left|\int_0^T(f_n)_s\diff X_s\right|^2]=\mathbb{E}[\left|c_n\right|^2\left|(X_{\tau\wedge T}-X_{\sigma\wedge T})\right|^2]\). If this is to be UI, then in particular, it is bounded in \(L^1\). But \(\sup_{n\in\mathbb{N}}\mathbb{E}[c_n^2(X_{\tau\wedge T}-X_{\sigma\wedge T})^2]=(\sup_{n\in\mathbb{N}}c_n^2)\mathbb{E}[(X_{\tau\wedge T}-X_{\sigma\wedge T})^2]<\infty\) by assumption. But \(\sup_{n\in\mathbb{N}}c_n^2=\infty\) hence we must have \(\mathbb{E}[(X_{\tau\wedge T}-X_{\sigma\wedge T})^2]=0\) which implies \(X_{\tau\wedge T}-X_{\sigma\wedge T}=0\quad\mathbb{P}\)-a.s. As this holds for all \(\mathbb{F}\)-stopping times \(\tau,\sigma\), in particular deterministic ones, we have \(X\equiv 0\) almost surely. Hence, there can't be a (non-trivial) condition on \(X\) making the family \(\left\{\left|\int_{0}^T\vartheta_s^n\diff X_s\right|^2\right\}_{n\in\mathbb{N}}\) UI, as the latter being UI for any sequence \((\vartheta^n)_{n\in\mathbb{N}}\) implies that \(X\equiv 0\).
\bigbreak
\textcolor{blue}{As an initial observation, recall that \(\int_{0}^T\vartheta_s^n\diff X_s\xrightarrow{\mathbb{P}}\int_{0}^T\vartheta^*_s\diff X_s\;\) implies that there exists a subsequence \(\left\{\int_{0}^T\vartheta_s^{n(k)}\diff X_s\right\}_{k\in\mathbb{N}}\) that converges \(\mathbb{P}\)-a.s. to \(\int_{0}^T\vartheta_s^{*}\diff X_s\). Hence, on a full measure set, it is (uniformly) bounded as a real converging sequence (assuming that \(\int_{0}^T\vartheta_s^{*}\diff X_s<\infty\) on that full measure set). If it is only \(\mathbb{P}\)-a.s. finite, take the intersection of the respective full measure sets to reach the same conclusion.) Hence, there exists a positive constant \(M\in\mathbb{R}_+\) such that \(\left|\int_{0}^T\vartheta_s^{n(k)}\diff X_s\right|^2<M^2\quad \mathbb{P}\)-a.s. for all \(k\in\mathbb{N}\). But \(M^2\) is integrable; hence, the family \(\left\{\left|\int_{0}^T\vartheta_s^{n(k)}\diff X_s\right|^2\right\}_{k\in\mathbb{N}}\) is uniformly integrable.\footnote{If \(H\geq |X_n|\) for all \(n\in\mathbb{N}\), we have \(\mathbb{E}[|X|\mathbbm{1}_{|X|\geq K}]\leq\mathbb{E}[H\mathbbm{1}_{H\geq K}]\) for any \(K>0\). As the left hand side going to 0 as \(K\to\infty\) it gives us the uniform integrability of \(\{X_n\}_{n\in\mathbb{N}}\)}. Thus, we get }
\\\\
\textcolor{red}{This conclusion is WRONG. Notice that the bound \(M = M(\omega)\) and hence doesn't define a random variable \(H\in L^2\)(at least not because it is constant almost surely). Indeed the argument used, goes as follows: For every \(\omega\) the sequence \(\left(\int_{0}^T\vartheta_s^{n(k)}\diff X_s\right)(\omega)\) is bounded as a real converging sequence, hence \(\left|\left(\int_{0}^T\vartheta_s^{n(k)}\diff X_s\right)(\omega)\right|\leq M(\omega)\).}
\\\\
\begin{equation}
\int_{0}^T\vartheta_s^{n(k)}\diff X_s \xrightarrow{L^2(\mathbb{P})} \int_{0}^T\vartheta^*_s\diff X_s\;\text{ as } k\to\infty
\end{equation}
Hence, restricting ourselves to this subsequence (which obviously still converges in probability to the same limit), we have 
\begin{align*}
    &\inf_{k\in\mathbb{N}}\Vert[H-V_0^*-\int_{0}^T\vartheta_s^{n(k)}\diff X_s]\rVert_{L^2}\leq \liminf_{k\to\infty}\lVert H-V_0^*-\int_{0}^T\vartheta_s^{n(k)}\diff X_s]\rVert_{L^2}\\&\leq \liminf_{k\to\infty}\lVert H-V_0^*-\int_{0}^T\vartheta_s^{*}\diff X_s\rVert_{L^2}+\lVert \int_{0}^T\vartheta_s^{*}\diff X_s -\int_{0}^T\vartheta_s^{n(k)}\diff X_s\rVert_{L^2}
    \\&=\lVert H-V_0^*-\int_{0}^T\vartheta_s^{*}\diff X_s\rVert_{L^2}
\end{align*}
where \(\lVert H-V_0^*-\int_{0}^T\vartheta_s^{*}\diff X_s\rVert_{L^2}=\inf_{\vartheta\in\Theta}\lVert H-V_0^*-\int_{0}^T\vartheta_s\diff X_s\rVert_{L^2}\).
By \cref{rem:inf and sequences}, the existence of this sequence is equivalent to a more general statement:
\begin{theorem}[Semimartingale Deep Hedging]\label{thm:general deep hedging}
    Let \(X\) be a semimartingale and \(\Theta\subset\Theta_2\). Let \((V_0^*,\vartheta^*)\) be the solution to \cref{thm:Hilbert Projection for Mean Sqaure Hedging}. Then there exists \((\vartheta_n)_{n\in\mathbb{N}}\subset\mathcal{S}(\mathcal{\psi})\) such that:
    \begin{equation}
        \min_{\vartheta\in \Theta}\lVert H-V_0^*-\int_{0}^T\vartheta_s\diff X_s\rVert_{L^2}=\inf_{\vartheta\in\mathcal{S}(\psi)}\lVert H-V_0^*-\int_{0}^T\vartheta_s\diff X_s\rVert_{L^2}\;,
    \end{equation}
\end{theorem}
where \(\lVert H-V_0^*-\int_{0}^T\vartheta^*_s\diff X_s\rVert_{L^2}=\min_{\vartheta\in \Theta}\lVert H-V_0^*-\int_{0}^T\vartheta_s\diff X_s\rVert_{L^2}\).
\\\\
\newpage
\underline{Possible strategies}:
In order to get deep-hedging result similar to \cref{thm:general deep hedging} we need 
\begin{equation}
    \int_0^T\vartheta^n_s\diff X_s \overset{L^2}{\to} \int_0^T\vartheta_s\diff X_s
\end{equation}
For an approximating sequence \((\vartheta_n)_{n\in\mathbb{N}}\subset\mathcal{S}(\psi)\) of \(\vartheta\in\tilde{\Theta}\), where the sense of approximation is to be determinded and the set \(\tilde{\Theta}\) such that \(\mathcal{S}(\psi)\subset\tilde{\Theta}\). Assuming this is done, we have
\begin{align*}
  &\inf_{\vartheta\in\mathcal{S}(\psi)}\lVert H-V_0^*-\int_0^T \vartheta_s^n\diff X_s\rVert_{L^2}\\&\leq
    \inf_{n\in\mathbb{N}}\lVert H-V_0^*-\int_0^T \vartheta_s^n\diff X_s\rVert_{L^2}\leq \liminf_{n\to\infty}\lVert H-V_0^*-\int_0^T \vartheta_s^n\diff X_s\rVert_{L^2}
    \\&\leq \liminf_{n\to\infty}\lVert H-V_0^*-\int_0^T \vartheta_s^*\diff X_s\rVert_{L^2}+\lVert \int_0^T \vartheta_s^n\diff X_s-\int_0^T \vartheta_s^*\diff X_s\rVert_{L^2}\\&=\lVert H-V_0^*-\int_0^T \vartheta_s^*\diff X_s\rVert_{L^2}\;,
\end{align*}
where \(\lVert H-V_0^*-\int_0^T \vartheta_s^*\diff X_s\rVert_{L^2}=\argmin_{\vartheta\in\Theta}\lVert H-V_0^*-\int_0^T \vartheta_s^*\diff X_s\rVert_{L^2}\) for some set \(\Theta\subset\Theta_2\). On the other hand we have \(\inf_{\vartheta\in\tilde{\Theta}}\lVert H-V_0^*-\int_0^T \vartheta_s^n\diff X_s\rVert_{L^2}\leq\inf_{\vartheta\in\mathcal{S}(\psi)}\lVert H-V_0^*-\int_0^T \vartheta_s^n\diff X_s\rVert_{L^2}\), because \(\mathcal{S}(\psi)\subset\tilde{\Theta}\). Thus we get equality. \\
We consider two possible strategies.
\bigbreak
(1): From \parencite[Thm~5.29.]{Arandjelovic2025} we have \(\mathcal{S}(\psi)\) is dense in \(\mathcal{S}\) which is dense in \(\mathcal{L}\), the space of càglàd processes (see Protter for this). Then for \(\vartheta_n\overset{\mathbb{P}-ucp}{\to}\vartheta\) where \(\vartheta\in \mathcal{L}\), we have 
\begin{align*}
    \int \vartheta_s^n\diff X_s\overset{\mathbb{P}-ucp}{\to}\int \vartheta_s\diff X_s\implies\int_0^T \vartheta_s^n\diff X_s\overset{\mathbb{P}}{\to}\int_0^T\vartheta_s\diff X_s
\end{align*}
and 
\begin{align*}
    &\left(\vartheta_s^n\diff X_s\overset{\mathbb{P}}{\to}\int_0^T\vartheta_s\diff X_s \right) + \left(\left\{\left|\int_{0}^T\vartheta_s^{n}\diff X_s\right|^2\right\}_{n\in\mathbb{N}}\text{uniformly integrable}\right)\\&\implies \left(\int_0^T\vartheta_s^n\diff X_s\overset{L^2}{\to}\int_0^T\vartheta_s\diff X_s \right)
\end{align*}
WARNING: Notice that for sequences of the form \(\vartheta_n=c_n\cdot\mathbbm{1}_{[\tau,\sigma)}\in\mathcal{S}(\psi)\) for \(\lim_{n\to\infty}c_n=\infty\), requiring that \(\left\{\left|\int_{0}^T\vartheta_s^{n}\diff X_s\right|^2\right\}_{n\in\mathbb{N}}\) is uniformly integrable implies \(X=0\) as a process. This is a pathological example as it doesn't converge to an interesting limiting process \(\vartheta\), which we assume. This just shows that  \(\left\{\left|\int_{0}^T\vartheta_s^{n}\diff X_s\right|^2\right\}_{n\in\mathbb{N}}\) can't be UI for all \((\vartheta_n)_{n\in\mathbb{N}}\subset\mathcal{S}(\psi)\). 
\bigbreak
(2): Notice that \(Y_n\overset{L^2-ucp}{\to}Y\) i.e. \(\lim_{n\to\infty}\mathbb{E}[\sup_{0\leq t\leq T}|Y_t^n-Y_t|^2]=0\) implies \(Y^n_T\overset{L^2}{\to}Y_T\). Hence, we want a version of \parencite[Thm~5.29]{Arandjelovic2025} with \(L^2-ucp\) convergence of the integrals. For this, we might want to use the fact that for \(\Phi(\cdot)=(\cdot)^2\), \parencite[Thm~5.28]{Arandjelovic2025} implies that \(\mathcal{S}(\psi)\) is \(L^2-ucp\) dense in \(\mathcal{S}\) and use some kind of \(L^2-ucp\) continuity of the stochastic integral \(H\to H\bullet X\) for well chosen classes for \(H\) and \(X\). An clearly weaker assumption would be to ask for convergence in \(\mathcal{H}^2\) which implies the \(L^2\) convergence of the end-points (Doob). A well known case where we have the continuity of the integral map w.r.t. to the \(L^2\)-convergence of the end-points is the isometry property (Itô isomtrey) of the map 
\begin{align}
    \begin{cases}
        H\to H\bullet M\\
        L^2(M)\to\mathcal{H}^2_0
    \end{cases}
\end{align}
for \(M\in\mathcal{H}^2_{0,loc}\)
Indeed in this case, by \parencite[Thm~5.32]{Arandjelovic2025} we even have that \(\mathcal{E}(\psi)\) is dense in \(L^p(M)\) for \(p\in[1,\infty)\) if \(M\in\mathcal{H}_0^{p+\epsilon}\) for some \(\epsilon>0\). (Then we also have that \(\mathcal{S}(\psi)\) is dense in \(L^p(M)\)). Hence we have 
\begin{equation}
    \lVert H\bullet M\rVert_{\mathcal{H}^2}=\lVert H\rVert_{L^2(M)}
\end{equation}
Using the sequence \(H_n=\vartheta_n-\vartheta\) and the fact that convergence in \(\mathcal{H}^2\) implies convergence in \(L^2-ucp\) (by Doob's maximal inequality, which is applicable because \(H\bullet M\) is a square integrable martingale in this case) we recover \parencite[Thm~5.35]{Arandjelovic2025}.
\\\\
To get a slight extension of \parencite[Thm~5.35]{Arandjelovic2025}, we notice that we don't need the full isometry property; rather, we need the boundedness (continuity) of the map \(H\to H \bullet M\) and the corresponding density of \(\mathcal{S}(\psi)\). Hence, for a semimartingale 
\begin{equation}
    X = X_0 + A + M
\end{equation}
for \(A\) a finite-variation process and \(M\in\mathcal{H}^2_{0,loc}\) (the addition of \(A\) destroys the Itô-isometry property) we try to bound 
\begin{equation}
     \lVert \left(H\bullet X\right)_T\rVert_{L^2}\leq K_X\lVert H\rVert_{L^2(M)}
\end{equation}
for some constant \(K_X = K(X)\). By \(|a+b|^2\leq 2|a|^2+2|b|^2\) it is enough to bound \(\lVert H\bullet A\rVert_{\mathcal{H}^2}\) and \(\lVert H\bullet M\rVert_{\mathcal{H}^2}\) both by \(\lVert H\rVert_{L^2(M)}\) and some multiplicative constant. Notice that for \(M\in\mathcal{H}^2_{0,loc}\) and \( H \in L^2(M)\), \(H\bullet M\) is a square integrable martingale; hence, we can bound \(\lVert H\bullet M\rVert_{\mathcal{H}^2}=\lVert (H\bullet M)_T\rVert_{L^2}\leq 4\lVert H\rVert_{L^2}\) by Doob. For the finite variation part, we might need assumptions on \(A\) (and hence on \(X\)).
\\
Based on this oberservations we consider the following approach. Consider a new space \(L^2(M,A)\) for a semimartingale \(X=X_0+M+A\) where \(A\) is a finite-variation process and \(M\in\mathcal{H}^2_{0,loc}\). Define it by \footnote{Note that in our case \([X]=[M]\) and \(L^2(M,A)\subset L^2(M)\), so \(\int_0^T\vartheta_s\diff M_s\in\mathcal{H}_0^2\); hence, \(\left(\int \vartheta_s\diff M_s\right)^2-[\int \vartheta_s\diff M_s]\) is a martingale starting at 0, and \([\int\vartheta_s\diff M_s]=\int \vartheta_s^2\diff [M]_s\). Hence, \(\mathbb{E}[\left(\int_0^T \vartheta_s\diff M_s\right)^2]=\mathbb{E}[\int_0^T \vartheta_s^2\diff [M]_s]\)}
\begin{align*}
H\in L^2(M,A) \iff \left(\mathbb{E}[\int_0^T H_s^2\diff [X]_s]+\mathbb{E}[\left(\int_0^T |H|_s\diff |A|_s\right)^2]\right)^{\frac{1}{2}}\coloneqq\lVert H\rVert_{L^2(M,A)}<\infty
\end{align*}\footnote{Inspired by the fact that for \(\int H_s\diff X_s = \int H_s\diff M_s+\int H_s\diff A_s\) and \(\mathbb{E}[\left(\int H_s\diff M_s+\int H_s\diff A_s\right)^2]\leq 2\mathbb{E}[\left(\int H_s \diff M_s\right)^2]+2\mathbb{E}[\left(\int H_s \diff A_s\right)^2]\leq2\cdot 4\mathbb{E}[\left(\int (H_s)^2 \diff [M]_s\right)^2]+2\mathbb{E}[\left(\int H_s \diff A_s\right)^2]\) from \(|a+b|^2\leq 2|a|^2+2|b|^2\) and BDG with \(C_2 = 4\).}\textcolor{blue}{One would still need to check that \(\lVert H\rVert_{L^2(M,A)}\) defines a norm}
where in our case \([X_0 + A + M]=[M]\). With this we have 
\begin{equation}
    \lVert\int_0^T\vartheta_s\diff X_s\rVert_{L^2}\leq K_X\lVert \vartheta \rVert_{L^2(M,A)}
\end{equation}
Indeed,
\begin{align*}\label{eq:bounding the L^2 norm of the stochastic integral}
    &\mathbb{E}[\left(\int_0^T H_s\diff X_s\right)^2]=\mathbb{E}[\left(\int_0^T H_s\diff M_s+\int_0^T H_s\diff A_s\right)^2]
    \\&
    \leq 2\mathbb{E}[\left(\int_0^T H_s \diff M_s\right)^2]+2\mathbb{E}[\left(\int_0^T H_s \diff A_s\right)^2]
    \\&
    \leq2\cdot 4\mathbb{E}[\int_0^T (H_s)^2 \diff [M]_s]+2\mathbb{E}[\left(\int_0^T H_s \diff A_s\right)^2]\leq
    \\&
     \leq2\cdot 4\mathbb{E}[\int_0^T (H_s)^2 \diff [M]_s]+2\mathbb{E}[\left(\int_0^T |H|_s \diff |A|_s\right)^2]\leq
    \lVert H\rVert_{L^2(M,A)}
\end{align*}
from \(|a+b|^2\leq 2|a|^2+2|b|^2\) and BDG with \(C_2 = 4\)\footnote{\(\mathbb{E}[\left(\int_0^TH_s\diff M_s\right)^2]\leq\mathbb{E}[\sup_{0\leq t\leq T}\left(\int_0^t H_s\diff M_s\right)^2]\leq C_2\mathbb{E}[\int_0^T(H_s)^2\diff [M]_s]\)}. Then we have
\begin{align}
    &\vartheta_n\overset{L^2(M,A)}{\to}\vartheta\implies\int_0^T\vartheta_s^n\diff X_s \overset{L^2}{\to}\int_0^T\vartheta_s\diff X_s
\end{align}
In this approach, we need to show that \(\lVert\cdot \rVert_{L^2(M,A)}\) is indeed a norm (or a semi-norm), that \(\mathcal{S}(\psi)\) is dense in \(L^2(M,A)\), and whether \(L^2(M,A)\) is rich enough to be interesting.
\section{Proof Attempt}
We present the first version of a generalization of \parencite[Thm~5.35]{Arandjelovic2025}. From \parencite[Thm.~5.32.]{Arandjelovic2025} we know that \(\mathcal{E}(\psi)\) is dense in \(L^p(M)\) for \(p\in[1,\infty)\) and \(M\in\mathcal{H}_0^{p+\epsilon}\) for some \(\epsilon>0\). 
Now consider a semimartingale \(X=X_0+M+A\), where \(M\in\mathcal{H}_0^{p+\epsilon}\) and \(A\) are adapted processes with RCLL trajectories, \(A_.(\omega)\) of finite variation and \(A_0=0\). Additionally, assume the following:
\begin{itemize}
    \item \(A<<\langle M\rangle\) path-wise (i.e. for every \(\omega\) as Lebesgue Stieltjes measures, with density process \(\lambda\).
    \item \(\operatorname*{ess\sup}\left(\langle M\rangle_{\infty}(\lambda^2+1)^*_{\infty}\right)<\infty\) where \((\lambda^2+1)^*_{\infty}=\sup_{t\geq 0}(\lambda^2+1)\).\footnote{Replace \(\infty\) by \(T\) in the finite horizon setting.}
\end{itemize}
Then we have, by density, of \(\mathcal{E}(\psi)\) in \(L^2(M)\) that
\begin{equation}
    \forall H\in L^2(M), \exists (H_n)_{n\in\mathbb{N}}\subset\mathcal{E}(\psi):\lim_{n\to\infty}\lVert H-H_n\rVert_{L^2(M)}=0\;.
\end{equation}
We can bound the \(L^2\) difference of the integrals by 
\begin{align}
    &\mathbb{E}[(H\bullet X)_{\infty}^2]\leq 2\left((\mathbb{E}[(H\bullet \langle M\rangle)_{\infty}^2]+\mathbb{E}[(H\bullet A)_{\infty}^2]\right)\\
    &\leq 2\left((\mathbb{E}[\langle M\rangle_{\infty}(H^2\bullet \langle M\rangle)_{\infty}]+\mathbb{E}[(H\lambda\bullet \langle M\rangle)_{\infty}^2]\right)\\
    &\leq2\left((\mathbb{E}[\langle M\rangle_{\infty}(H^2\bullet \langle M\rangle)_{\infty}]+\mathbb{E}[\langle M\rangle_{\infty}(H^2\lambda^2\bullet \langle M\rangle)_{\infty}]\right)\\
    &=2\mathbb{E}[\langle M\rangle_{\infty}(H^2(1+\lambda^2)\bullet\langle M\rangle)_{\infty}]\\
    &\leq 2\mathbb{E}[\langle M\rangle_{\infty}(1+\lambda^2)^{*}_{\infty}(H^2\bullet\langle M\rangle)_{\infty}]\\
    &\leq2\operatorname*{ess\sup}(\langle M\rangle_{\infty}(1+\lambda^2)^{*}_{\infty})\cdot\mathbb{E}[H^2\bullet \langle M\rangle]\\
    &=C(X)\lVert H\rVert_{L^2(M)}^2\;.
\end{align}
We used Cauchy Schwarz in step 2 and 3, along with the absolute continuity of \(A<<\langle M\rangle\). In step 5, we used the Hölder inequality with \(p=1,q=\infty\).\\
By assumption \(C(X)<\infty\), hence substituting for \(H\), \(Z^n=H-H^n\), we get that 
\begin{equation}
    \lim_{n\to\infty}\mathbb{E}[(Z^n\bullet X)_{\infty}^2]\leq\lim_{n\to\infty}C(X)\lVert H\rVert_{L^2(M)}^2=0
\end{equation}
By the argument given earlier, this allows for a deep mean-variance hedging result for a subclass of semimartingales.
\section{Different Approach}
We present a slightly different approach based on the observation \cref{eq:bounding the L^2 norm of the stochastic integral}. Notice that under the assumption that \(A<<\langle M\rangle\) path-wise (i.e. for every \(\omega\) as Lebesgue Stieltjes measures, with density process \(\lambda\) we have:
\begin{align*}
&\lVert H\rVert_{L^2(M,A)}^2=\mathbb{E}[\int_0^\infty (H_s)^2\diff [M]_s]+\mathbb{E}[\left(\int_0^\infty |H_s||\lambda_s|\diff[M]_s\right)^2]\\&\leq\mathbb{E}[\int_0^\infty (H_s)^2\diff [M]_s+\mathbb{E}[\int_0^{\infty}(H_s)^2\diff [M]_s\cdot\int_0^{\infty}(\lambda_s)^2\diff [M]_s]
\\&
\leq\mathbb{E}[\int_0^\infty (H_s)^2\diff [M]_s+\mathbb{E}[\left(\int_0^{\infty}(H_s)^2\diff [M]_s\right)^{p}]^{1/p}\cdot\mathbb{E}[\left(\int_0^{\infty}(\lambda_s)^2\diff [M]_s\right)^{q}]^{1/q}
\;,
\end{align*}
where we used the fact that if \(A<<[M]\) with density \(\lambda\) then \(|A|<<[M]\) with density \(|\lambda|\). Furthermore we used the Kunita-Watanabe inequality in the second step and Hölder's inequality in the third step. In the case \(p=1,q=\infty\) the above reads
\begin{equation}
    \lVert H\rVert_{L^2(M,A)}\leq\lVert H\rVert_{L^2(M)}\left(1+\operatorname*{ess\sup}\left(\int_0^{\infty}(\lambda_s)^2\diff [M]_s\right)\right)\;.
\end{equation}
Notice \(\int_0^{\infty}(\lambda_s)^2\diff [M]_s=\int_0^{\infty}|\lambda_s|\diff A_s\leq \int_0^{\infty}|\lambda_s|\diff |A|_s\). We impose the additional assumption
\begin{equation}
    \operatorname*{ess\sup}\int_0^{\infty}|\lambda_s|\diff |A|_s<\infty\;.
\end{equation}
Then, because we have the trivial inequality
\begin{align}\label{eq:essential boundedness assumption}
    \mathbb{E}[\int_0^{\infty}H_s^2\diff[M]_s]\leq  \mathbb{E}[\int_0^{\infty}H_s^2\diff[M]_s]+  \mathbb{E}[\left(\int_0^{\infty}|H_s||\lambda_s|\diff[M]_s\right)^2]=\lVert H\rVert_{L^2(M,A)}
\end{align}
we have that the norms on \(L^2(M)\) and \(L^2(M,A)\),under the condition that \(X = X_0+M+A\) with \(A<<[M]\) and \(M\in\mathcal{H}^{p+\epsilon}_0\) for some \(p\in[0,\infty)\) and \(\epsilon>0\), along with the additional assumption \cref{eq:essential boundedness assumption}, are equivalent. Because both spaces \(L^2(M)\) and \(L^2(M,A)\) are subspaces of the space of \(\mathcal
P\)-measurable functions in, where \(\mathcal{P}\) is the predictable \(\sigma\)-algebra on \(\Omega\times[0,\infty)\), they are defined through their norms; the spaces are equal:
\begin{equation}
    L^2(M) =L^2(M,A)\;.
\end{equation}
Hence, density of \(\mathcal{E}(\psi)\) in \(L^2(M)\) is obviously equivalent to density of \(\mathcal{E}(\psi)\) in \(L^2(M,A)\) under the stated assumptions on \(M\) and \(X\). Hence, for every \(H\in L^2(M,A)\) there exists \((H_n)_{n\in\mathbb{N}}\subset \mathcal{E}(\psi)\) such that \(\lim_{n\to\infty}\lVert H-H_n\rVert_{L^2(M,A)}=0\). Then we have, by the observation \cref{eq:bounding the L^2 norm of the stochastic integral}, with \(T=\infty\) and \(H_s\coloneqq (H-H_n)_s\)
\begin{equation}
    \lVert (H-H^n)\bullet X\rVert_{L^2}
    \lesssim\lVert H-H^n\rVert_{L^2(M,A)}\;,
\end{equation}
thus
\begin{equation}
    \lim_{n\to\infty}\lVert (H-H^n)\bullet X\rVert_{L^2}=0
\end{equation}
\bigbreak
The two approaches differ in their additional assumption
\begin{enumerate}
    \item \(\operatorname*{ess\sup}(\langle M\rangle_{\infty}(1+\lambda^2)^{*}_{\infty})<\infty\)
    \item  \(\operatorname*{ess\sup}\int_0^{\infty}|\lambda_s|\diff |A|_s<\infty\)
\end{enumerate}
Notice that
\begin{equation}
    \int_0^\infty(\lambda_s)^2\diff [M]_s\leq\int_0^{\infty}(1+\lambda_s^2)\diff[M]_s\leq(1+\lambda_s^2)_{\infty}^*[M]_{\infty}\,
\end{equation}
where we used the non-negativity and monotonicity of the process \([M]\). Hence, (2.) is a weaker assumption than (1.).
\textcolor{blue}{potential problem with this approach. The norms are equivalent but not really norms. First one needs to check that \(\lVert \cdot\rVert_{L^2(M,A)}\) defines at least a subadditiv, homogenous map and then that the process which are identified by \(\lVert X-Y\rVert_{L^2(M,A)}=0\) are also identified in \(L^2(M)\), that is \(\lVert X-Y\rVert_{L^2(M)}=0\). So that the norms are norms on the same quotient space under some equivalence relation identifying processes}
First, notice that \(\lVert H\rVert_{L^2(M,A)}=0\) is subadditive
Under the assumption that \(A<<[M]\) has density \(\lambda\), we have that \(\lVert X-Y\rVert_{L^2(M)}=0\implies\lVert X-Y\rVert_{L^2(M,A)}=0\). This follows from the fact that \(c\lVert X\rVert_{L^2(M)}\leq\lVert X\rVert_{L^2(M,A)}\leq C\lVert X\rVert_{L^2(M)}\) for \(c=1,C=(1+\operatorname*{ess\sup}\int_0^{\infty}|\lambda_s|\diff |A|_s)\). By assumption, \(c,C\) are finite. This implies that under the equivalence relation \(X\sim Y\) for \(X\sim Y\) for \(\lVert X-Y\rVert_{L^2(M,A)}=0\). Hence, they are norms on the same space \(L^2(\Omega\times [0,T],\mathcal{P})/\sim\), and in particular, it makes sense to say that they are equivalent. 
\bigbreak
This proves the following result:
\begin{theorem}[Semimartingale Deep Hedging]\label{thm:general deep hedging}
    Let \(X=X_0+M+A\) be a semimartingale with \(M\in \mathcal{H}_{0}^{2+\epsilon}\) for some \(\epsilon>0\), \(A<<[M]\) with density \(\lambda=(\lambda_s)_{s\in\mathbb{R}_+}\) and with \(\operatorname*{ess\sup}\int_0^{\infty}|\lambda_s|\diff |A|_s<\infty\) \. Let \(H\in L^2(\mathbb{P})\) and \((V_0^*,\vartheta^*)\) be the solution to \cref{thm:Hilbert Projection for Mean Sqaure Hedging}. Then there exists \((\vartheta_n)_{n\in\mathbb{N}}\subset\mathcal{E}(\mathcal{\psi})\) such that:
    \begin{equation}
        \min_{\vartheta\in \Theta}\lVert H-V_0^*-\int_{0}^T\vartheta_s\diff X_s\rVert_{L^2}=\inf_{\vartheta\in\mathcal{E}(\psi)}\lVert H-V_0^*-\int_{0}^T\vartheta_s\diff X_s\rVert_{L^2}\;,
    \end{equation}
\end{theorem}
where \(\lVert H-V_0^*-\int_{0}^T\vartheta^*_s\diff X_s\rVert_{L^2}=\min_{\vartheta\in \Theta}\lVert H-V_0^*-\int_{0}^T\vartheta_s\diff X_s\rVert_{L^2}\).
\begin{proof}
    The proof is given above and relies on the fact that in this case \(L^2(M)=L^2(M,A)\). By \parencite[Thm.~5.32]{Arandjelovic2025} \(\mathcal{E}(\psi)\) is dense in \(L^2(M)\) and hence also trivially in \(L^2(M,A)\) as well. The only difference is that the optimal \(V_0^*,H^*\) are for the more general problem with \(X= X_0+M+A\). The proof is the same as in \parencite[THm.~5.35]{Arandjelovic2025} as for any \(\delta>0\) there exists \(U\in\mathcal{E}(\psi)\) such that \(\lVert(U^*\bullet X)_{\infty}-(U\bullet X)_{\infty}\rVert_{L^2}\leq \lVert U^*-U\rVert_{L^2(M,A)}<\delta\) by the calculation in \cref{eq:bounding the L^2 norm of the stochastic integral}.
\end{proof}
One should give examples, if they exist, of finite variation processes \(A\) satisfying the assmuptions:
\begin{itemize}
    \item \(A<<[M]\) with density \(\lambda\)
    \item \(\operatorname*{ess\sup}\int_0^{\infty}|\lambda_s|\diff |A|_s<\infty\)
\end{itemize}

Furthermore we will think of a more general approach generalizing \parencite[Thm.~5.35]{Arandjelovic2025} again with \(X\) a semimartingale and approximants from \(\mathcal{S}(\psi)\). For this we might think again about strategy (1) and if there are sensible conditions on \(\vartheta\) that transfer to conditions on \(\vartheta_n\) such that the respective family of stochastic integrals (more precisely the square processes) become a uniformly integrable family.

\textcolor{blue}{Possible, non-trivial process satisfying assumptions: Let}
\begin{equation}
    S_t = \lambda B_t + \mu t
\end{equation}
Then \(\langle S\rangle_t= \lambda^2 t\) and \(A_t=\mu t\).
\end{document}
% so it automatically matches the exact same layout/preamble.

\section{Proof attempts}

% Put your drafts here. Example skeleton:
%
%\begin{proof}[Proof attempt (Lemma~\ref{lem:...})]
%    ...
%\end{proof}
\section{Extension to semi-martingale case}
\cref{thm:cont-deep-ms-hedging-martingale} was in the realm of classical martingale stochastic calculus, with integrators \(M\in\mathcal{H}^{2}\)(even \(\mathcal{H}^{2+\epsilon}\) and integrands \(\vartheta\in L^2(M)\). In light of \cref{thm:Approx of stoch. Integrals}, we aim to prove a similar result as in \cref{thm:cont-deep-ms-hedging-martingale} for the more general case where \(M\) is a semimartingale. Following \parencite{Schweizer2001GuidedTourQuadraticHedging} we restrict the class of possible strategies to \(\vartheta_{GLP}\) such that \(G_T(\vartheta_{GLP})\) is a closed subspace in \(L^2(\mathbb{P})\). In this case, we know that the minimization problem has a solution, as \(\mathcal{A}=\mathbb{R}+G_T(\Theta_{GLP})\) is a closed subspace in \(L^2(\mathbb{P})\).; this means there exists an optimal trading strategy \((V_0^*,\vartheta^*)\in\mathcal{A}=\mathbb{R}+G_T(\Theta_{GLP})\) (\textcolor{blue}{I guess we don't even need strong conditions to get \(\mathcal{A}\) closed as we don't need explicit expression of the optimal trading strategy at first, hence we can directly project on \(\bar{\mathcal{A}}\) which is closed in \(L^2(\mathbb{P})\). In that line of thought, any \(\Theta\subset\Theta_2\) is sufficient for the existence of a solution in \(\bar{\mathcal{A}}\)}.) Then we have the following consequence of the Hilbert Projection theorem
\begin{theorem}[Hilbert Projection for Mean Sqaure Hedging]\label{thm:Hilbert Projection for Mean Sqaure Hedging}
    Let \(\Theta\subset\Theta_2\) and \(\mathcal{A}=\mathbb{R}+G_T(\Theta)\subset L^2(\mathbb{P})\). Then the \(L^2\)-closure \(\bar{\mathcal{A}}\) is a closed subspace of \(L^2(\mathbb{P})\). Hence, there exists \((V_0^*,\vartheta^*)\) such that 
    \begin{equation}
   \lVert H-V_0^*-\int_{0}^T\vartheta^*_s\diff X_s\rVert_{L^2}= \inf_{(c,\vartheta)\in\bar{\mathcal{A}}}\lVert H-c-\int_{0}^T\vartheta_s\diff X_s\rVert_{L^2}
    \end{equation}
    and the solution is unique. 
\end{theorem}
\begin{remark}
    Notice, we can study the related problem where we fix \(V_0\) (corresponds to the situation where the price \(V_0\) of the liability \(H\) might be determined by the market) and study
    \begin{equation}
        \inf_{\vartheta\in\Theta}\lVert H-V_0-\int_{0}^T\vartheta_s\diff X_s\rVert_{L^2}
    \end{equation}
    Again we can consider the the \(L^2\) closure of \(G_T(\Theta)=\left\{\int_0^T \vartheta_s\diff X_s:\vartheta\in\Theta\right\}\) and use the Hilbert Projection theorem on this closed subspace of \(L^2\).
\end{remark}
Then, crucially, notice that from \cref{thm:Approx of stoch. Integrals} we obtain the convergence in probability of the final values of the respective integral processes. Indeed, the \(\mathbb{P}\)-\(\mathbb{P}\)-ucp convergence of the integral processes, restricted to a finite time horizon \(T\), reads as:
\begin{equation}
    \sup_{0\leq t\leq T}\left|\int_{0}^t\vartheta_s^n\diff X_s-\int_{0}^t\vartheta^*_s\diff X_s\right|\to 0\quad\text{in Probability}
\end{equation}
where \(\{\vartheta^n\}_{n\in\mathbb{N}}\in\mathcal{S}(\psi)\) is a minimizing sequence from \cref{thm:Approx of stoch. Integrals}
\\\\
In particular, this implies 
\begin{equation}
    \int_{0}^T\vartheta_s^n\diff X_s\xrightarrow{\mathbb{P}}\int_{0}^T\vartheta^*_s\diff X_s\;,
\end{equation}
where convergence holds in probability. Imposing the additional assumption that \(\left\{\left|\int_{0}^T\vartheta_s^n\diff X_s\right|^2\right\}_{n\in\mathbb{N}}\) is a uniformly integrable family of random variables  leads to the \(L^2\) convergence\footnote{In particular, the condition also implies that the random variable \(\int_{0}^T\vartheta_s^n\diff X_s\in L^2(\mathbb{P)}\quad\forall n\in\mathbb{N}\)}: 
\begin{equation}
\int_{0}^T\vartheta_s^n\diff X_s \xrightarrow{L^2(\mathbb{P})} \int_{0}^T\vartheta^*_s\diff X_s\;.
\end{equation}
We do NOT want the uniform integrability of \(\left\{\left|\int_{0}^T\vartheta_s^n\diff X_s\right|^2\right\}_{n\in\mathbb{N}}\) to follow from conditions on a sequence of trading strategies. Indeed, \cref{thm:Approx of stoch. Integrals} asserts the existence of a sequence \(\{\vartheta^n\}_{n\in\mathbb{N}}\subset\mathcal{S}(\psi)\). We don't want a theorem that says something like: 
\\\\
"If that sequence happens to be square uniformly integrable, then it follows...."
\\\\
Hence, if necessary, we shall impose conditions only on the integrator \(X\). Recall that for \(\vartheta_n\in\mathcal{S}(\psi)\) that is \(\vartheta^n_t=\varphi_0\mathbbm{1}_{\{0\}}+\sum_{i=1}^l\varphi_i^n\mathbbm{1}_{\tau_i,\tau_{i+1}}\) and \(\varphi_i^{(n)}=f^{(n)}(Y^{\tau_i}_{t_1},\dots,Y^{\tau_i}_{t_r})\) where \(f\in\mathcal{NN}((\mathbb{R}^{d\times r})^*,\psi)\)\footnote{the parameter \(r\in\mathbb{N}\) is not further specified. It relies to the existence result of \parencite[Thm.~5.23.]{Arandjelovic2025}}, the stochastic integral takes the form \(\int_0^T\vartheta^n_s\diff X_s = \sum_{i=1}^k \varphi^{(n)}_i(X_{\tau_i\wedge T}-X_{\tau_{i+1}\wedge T})\). Note that \(\left|\int_0^T\vartheta^n_s\diff X_s\right|^2 \leq \sum_{i=1}^k \left|\varphi^{(n)}_i\right|^2\left|(X_{\tau_i\wedge T}-X_{\tau_{i+1}\wedge T})\right|^2\). But there can't exist a uniform bound (in n) on \(|\varphi_i^n|\) as this would imply that neural networks can't approximate constants above this bound which is absurd. Also, notice for \(\tau<\sigma\) two stopping times, that \(f_n = c_n\mathbbm{1}_{[\tau,\sigma)}\in\mathcal{S}(\psi)\) for all \(c_n\in\mathbb{R}\). Take \(c_n\to\infty\) as \(n\to\infty\) a real deterministic sequence, then \(\int_0^T(f_n)_s\diff X_s=c_n(X_{\tau\wedge T}-X_{\sigma\wedge T})\), and hence \(\mathbb{E}[\left|\int_0^T(f_n)_s\diff X_s\right|^2]=\mathbb{E}[\left|c_n\right|^2\left|(X_{\tau\wedge T}-X_{\sigma\wedge T})\right|^2]\). If this is to be UI, then in particular, it is bounded in \(L^1\). But \(\sup_{n\in\mathbb{N}}\mathbb{E}[c_n^2(X_{\tau\wedge T}-X_{\sigma\wedge T})^2]=(\sup_{n\in\mathbb{N}}c_n^2)\mathbb{E}[(X_{\tau\wedge T}-X_{\sigma\wedge T})^2]<\infty\) by assumption. But \(\sup_{n\in\mathbb{N}}c_n^2=\infty\) hence we must have \(\mathbb{E}[(X_{\tau\wedge T}-X_{\sigma\wedge T})^2]=0\) which implies \(X_{\tau\wedge T}-X_{\sigma\wedge T}=0\quad\mathbb{P}\)-a.s. As this holds for all \(\mathbb{F}\)-stopping times \(\tau,\sigma\), in particular deterministic ones, we have \(X\equiv 0\) almost surely. Hence, there can't be a (non-trivial) condition on \(X\) making the family \(\left\{\left|\int_{0}^T\vartheta_s^n\diff X_s\right|^2\right\}_{n\in\mathbb{N}}\) UI, as the latter being UI for any sequence \((\vartheta^n)_{n\in\mathbb{N}}\) implies that \(X\equiv 0\).
\bigbreak
\textcolor{blue}{As an initial observation, recall that \(\int_{0}^T\vartheta_s^n\diff X_s\xrightarrow{\mathbb{P}}\int_{0}^T\vartheta^*_s\diff X_s\;\) implies that there exists a subsequence \(\left\{\int_{0}^T\vartheta_s^{n(k)}\diff X_s\right\}_{k\in\mathbb{N}}\) that converges \(\mathbb{P}\)-a.s. to \(\int_{0}^T\vartheta_s^{*}\diff X_s\). Hence, on a full measure set, it is (uniformly) bounded as a real converging sequence (assuming that \(\int_{0}^T\vartheta_s^{*}\diff X_s<\infty\) on that full measure set). If it is only \(\mathbb{P}\)-a.s. finite, take the intersection of the respective full measure sets to reach the same conclusion.) Hence, there exists a positive constant \(M\in\mathbb{R}_+\) such that \(\left|\int_{0}^T\vartheta_s^{n(k)}\diff X_s\right|^2<M^2\quad \mathbb{P}\)-a.s. for all \(k\in\mathbb{N}\). But \(M^2\) is integrable; hence, the family \(\left\{\left|\int_{0}^T\vartheta_s^{n(k)}\diff X_s\right|^2\right\}_{k\in\mathbb{N}}\) is uniformly integrable.\footnote{If \(H\geq |X_n|\) for all \(n\in\mathbb{N}\), we have \(\mathbb{E}[|X|\mathbbm{1}_{|X|\geq K}]\leq\mathbb{E}[H\mathbbm{1}_{H\geq K}]\) for any \(K>0\). As the left hand side going to 0 as \(K\to\infty\) it gives us the uniform integrability of \(\{X_n\}_{n\in\mathbb{N}}\)}. Thus, we get }
\\\\
\textcolor{red}{This conclusion is WRONG. Notice that the bound \(M = M(\omega)\) and hence doesn't define a random variable \(H\in L^2\)(at least not because it is constant almost surely). Indeed the argument used, goes as follows: For every \(\omega\) the sequence \(\left(\int_{0}^T\vartheta_s^{n(k)}\diff X_s\right)(\omega)\) is bounded as a real converging sequence, hence \(\left|\left(\int_{0}^T\vartheta_s^{n(k)}\diff X_s\right)(\omega)\right|\leq M(\omega)\).}
\\\\
\begin{equation}
\int_{0}^T\vartheta_s^{n(k)}\diff X_s \xrightarrow{L^2(\mathbb{P})} \int_{0}^T\vartheta^*_s\diff X_s\;\text{ as } k\to\infty
\end{equation}
Hence, restricting ourselves to this subsequence (which obviously still converges in probability to the same limit), we have 
\begin{align*}
    &\inf_{k\in\mathbb{N}}\Vert[H-V_0^*-\int_{0}^T\vartheta_s^{n(k)}\diff X_s]\rVert_{L^2}\leq \liminf_{k\to\infty}\lVert H-V_0^*-\int_{0}^T\vartheta_s^{n(k)}\diff X_s]\rVert_{L^2}\\&\leq \liminf_{k\to\infty}\lVert H-V_0^*-\int_{0}^T\vartheta_s^{*}\diff X_s\rVert_{L^2}+\lVert \int_{0}^T\vartheta_s^{*}\diff X_s -\int_{0}^T\vartheta_s^{n(k)}\diff X_s\rVert_{L^2}
    \\&=\lVert H-V_0^*-\int_{0}^T\vartheta_s^{*}\diff X_s\rVert_{L^2}
\end{align*}
where \(\lVert H-V_0^*-\int_{0}^T\vartheta_s^{*}\diff X_s\rVert_{L^2}=\inf_{\vartheta\in\Theta}\lVert H-V_0^*-\int_{0}^T\vartheta_s\diff X_s\rVert_{L^2}\).
By \cref{rem:inf and sequences}, the existence of this sequence is equivalent to a more general statement:
\begin{theorem}[Semimartingale Deep Hedging]\label{thm:general deep hedging}
    Let \(X\) be a semimartingale and \(\Theta\subset\Theta_2\). Let \((V_0^*,\vartheta^*)\) be the solution to \cref{thm:Hilbert Projection for Mean Sqaure Hedging}. Then there exists \((\vartheta_n)_{n\in\mathbb{N}}\subset\mathcal{S}(\mathcal{\psi})\) such that:
    \begin{equation}
        \min_{\vartheta\in \Theta}\lVert H-V_0^*-\int_{0}^T\vartheta_s\diff X_s\rVert_{L^2}=\inf_{\vartheta\in\mathcal{S}(\psi)}\lVert H-V_0^*-\int_{0}^T\vartheta_s\diff X_s\rVert_{L^2}\;,
    \end{equation}
\end{theorem}
where \(\lVert H-V_0^*-\int_{0}^T\vartheta^*_s\diff X_s\rVert_{L^2}=\min_{\vartheta\in \Theta}\lVert H-V_0^*-\int_{0}^T\vartheta_s\diff X_s\rVert_{L^2}\).
\\\\
\newpage
\underline{Possible strategies}:
In order to get deep-hedging result similar to \cref{thm:general deep hedging} we need 
\begin{equation}
    \int_0^T\vartheta^n_s\diff X_s \overset{L^2}{\to} \int_0^T\vartheta_s\diff X_s
\end{equation}
For an approximating sequence \((\vartheta_n)_{n\in\mathbb{N}}\subset\mathcal{S}(\psi)\) of \(\vartheta\in\tilde{\Theta}\), where the sense of approximation is to be determinded and the set \(\tilde{\Theta}\) such that \(\mathcal{S}(\psi)\subset\tilde{\Theta}\). Assuming this is done, we have
\begin{align*}
  &\inf_{\vartheta\in\mathcal{S}(\psi)}\lVert H-V_0^*-\int_0^T \vartheta_s^n\diff X_s\rVert_{L^2}\\&\leq
    \inf_{n\in\mathbb{N}}\lVert H-V_0^*-\int_0^T \vartheta_s^n\diff X_s\rVert_{L^2}\leq \liminf_{n\to\infty}\lVert H-V_0^*-\int_0^T \vartheta_s^n\diff X_s\rVert_{L^2}
    \\&\leq \liminf_{n\to\infty}\lVert H-V_0^*-\int_0^T \vartheta_s^*\diff X_s\rVert_{L^2}+\lVert \int_0^T \vartheta_s^n\diff X_s-\int_0^T \vartheta_s^*\diff X_s\rVert_{L^2}\\&=\lVert H-V_0^*-\int_0^T \vartheta_s^*\diff X_s\rVert_{L^2}\;,
\end{align*}
where \(\lVert H-V_0^*-\int_0^T \vartheta_s^*\diff X_s\rVert_{L^2}=\argmin_{\vartheta\in\Theta}\lVert H-V_0^*-\int_0^T \vartheta_s^*\diff X_s\rVert_{L^2}\) for some set \(\Theta\subset\Theta_2\). On the other hand we have \(\inf_{\vartheta\in\tilde{\Theta}}\lVert H-V_0^*-\int_0^T \vartheta_s^n\diff X_s\rVert_{L^2}\leq\inf_{\vartheta\in\mathcal{S}(\psi)}\lVert H-V_0^*-\int_0^T \vartheta_s^n\diff X_s\rVert_{L^2}\), because \(\mathcal{S}(\psi)\subset\tilde{\Theta}\). Thus we get equality. \\
We consider two possible strategies.
\bigbreak
(1): From \parencite[Thm~5.29.]{Arandjelovic2025} we have \(\mathcal{S}(\psi)\) is dense in \(\mathcal{S}\) which is dense in \(\mathcal{L}\), the space of càglàd processes (see Protter for this). Then for \(\vartheta_n\overset{\mathbb{P}-ucp}{\to}\vartheta\) where \(\vartheta\in \mathcal{L}\), we have 
\begin{align*}
    \int \vartheta_s^n\diff X_s\overset{\mathbb{P}-ucp}{\to}\int \vartheta_s\diff X_s\implies\int_0^T \vartheta_s^n\diff X_s\overset{\mathbb{P}}{\to}\int_0^T\vartheta_s\diff X_s
\end{align*}
and 
\begin{align*}
    &\left(\vartheta_s^n\diff X_s\overset{\mathbb{P}}{\to}\int_0^T\vartheta_s\diff X_s \right) + \left(\left\{\left|\int_{0}^T\vartheta_s^{n}\diff X_s\right|^2\right\}_{n\in\mathbb{N}}\text{uniformly integrable}\right)\\&\implies \left(\int_0^T\vartheta_s^n\diff X_s\overset{L^2}{\to}\int_0^T\vartheta_s\diff X_s \right)
\end{align*}
WARNING: Notice that for sequences of the form \(\vartheta_n=c_n\cdot\mathbbm{1}_{[\tau,\sigma)}\in\mathcal{S}(\psi)\) for \(\lim_{n\to\infty}c_n=\infty\), requiring that \(\left\{\left|\int_{0}^T\vartheta_s^{n}\diff X_s\right|^2\right\}_{n\in\mathbb{N}}\) is uniformly integrable implies \(X=0\) as a process. This is a pathological example as it doesn't converge to an interesting limiting process \(\vartheta\), which we assume. This just shows that  \(\left\{\left|\int_{0}^T\vartheta_s^{n}\diff X_s\right|^2\right\}_{n\in\mathbb{N}}\) can't be UI for all \((\vartheta_n)_{n\in\mathbb{N}}\subset\mathcal{S}(\psi)\). 
\bigbreak
(2): Notice that \(Y_n\overset{L^2-ucp}{\to}Y\) i.e. \(\lim_{n\to\infty}\mathbb{E}[\sup_{0\leq t\leq T}|Y_t^n-Y_t|^2]=0\) implies \(Y^n_T\overset{L^2}{\to}Y_T\). Hence, we want a version of \parencite[Thm~5.29]{Arandjelovic2025} with \(L^2-ucp\) convergence of the integrals. For this, we might want to use the fact that for \(\Phi(\cdot)=(\cdot)^2\), \parencite[Thm~5.28]{Arandjelovic2025} implies that \(\mathcal{S}(\psi)\) is \(L^2-ucp\) dense in \(\mathcal{S}\) and use some kind of \(L^2-ucp\) continuity of the stochastic integral \(H\to H\bullet X\) for well chosen classes for \(H\) and \(X\). An clearly weaker assumption would be to ask for convergence in \(\mathcal{H}^2\) which implies the \(L^2\) convergence of the end-points (Doob). A well known case where we have the continuity of the integral map w.r.t. to the \(L^2\)-convergence of the end-points is the isometry property (Itô isomtrey) of the map 
\begin{align}
    \begin{cases}
        H\to H\bullet M\\
        L^2(M)\to\mathcal{H}^2_0
    \end{cases}
\end{align}
for \(M\in\mathcal{H}^2_{0,loc}\)
Indeed in this case, by \parencite[Thm~5.32]{Arandjelovic2025} we even have that \(\mathcal{E}(\psi)\) is dense in \(L^p(M)\) for \(p\in[1,\infty)\) if \(M\in\mathcal{H}_0^{p+\epsilon}\) for some \(\epsilon>0\). (Then we also have that \(\mathcal{S}(\psi)\) is dense in \(L^p(M)\)). Hence we have 
\begin{equation}
    \lVert H\bullet M\rVert_{\mathcal{H}^2}=\lVert H\rVert_{L^2(M)}
\end{equation}
Using the sequence \(H_n=\vartheta_n-\vartheta\) and the fact that convergence in \(\mathcal{H}^2\) implies convergence in \(L^2-ucp\) (by Doob's maximal inequality, which is applicable because \(H\bullet M\) is a square integrable martingale in this case) we recover \parencite[Thm~5.35]{Arandjelovic2025}.
\\\\
To get a slight extension of \parencite[Thm~5.35]{Arandjelovic2025}, we notice that we don't need the full isometry property; rather, we need the boundedness (continuity) of the map \(H\to H \bullet M\) and the corresponding density of \(\mathcal{S}(\psi)\). Hence, for a semimartingale 
\begin{equation}
    X = X_0 + A + M
\end{equation}
for \(A\) a finite-variation process and \(M\in\mathcal{H}^2_{0,loc}\) (the addition of \(A\) destroys the Itô-isometry property) we try to bound 
\begin{equation}
     \lVert \left(H\bullet X\right)_T\rVert_{L^2}\leq K_X\lVert H\rVert_{L^2(M)}
\end{equation}
for some constant \(K_X = K(X)\). By \(|a+b|^2\leq 2|a|^2+2|b|^2\) it is enough to bound \(\lVert H\bullet A\rVert_{\mathcal{H}^2}\) and \(\lVert H\bullet M\rVert_{\mathcal{H}^2}\) both by \(\lVert H\rVert_{L^2(M)}\) and some multiplicative constant. Notice that for \(M\in\mathcal{H}^2_{0,loc}\) and \( H \in L^2(M)\), \(H\bullet M\) is a square integrable martingale; hence, we can bound \(\lVert H\bullet M\rVert_{\mathcal{H}^2}=\lVert (H\bullet M)_T\rVert_{L^2}\leq 4\lVert H\rVert_{L^2}\) by Doob. For the finite variation part, we might need assumptions on \(A\) (and hence on \(X\)).
\\
Based on this oberservations we consider the following approach. Consider a new space \(L^2(M,A)\) for a semimartingale \(X=X_0+M+A\) where \(A\) is a finite-variation process and \(M\in\mathcal{H}^2_{0,loc}\). Define it by \footnote{Note that in our case \([X]=[M]\) and \(L^2(M,A)\subset L^2(M)\), so \(\int_0^T\vartheta_s\diff M_s\in\mathcal{H}_0^2\); hence, \(\left(\int \vartheta_s\diff M_s\right)^2-[\int \vartheta_s\diff M_s]\) is a martingale starting at 0, and \([\int\vartheta_s\diff M_s]=\int \vartheta_s^2\diff [M]_s\). Hence, \(\mathbb{E}[\left(\int_0^T \vartheta_s\diff M_s\right)^2]=\mathbb{E}[\int_0^T \vartheta_s^2\diff [M]_s]\)}
\begin{align*}
H\in L^2(M,A) \iff \left(\mathbb{E}[\int_0^T H_s^2\diff [X]_s]+\mathbb{E}[\left(\int_0^T |H|_s\diff |A|_s\right)^2]\right)^{\frac{1}{2}}\coloneqq\lVert H\rVert_{L^2(M,A)}<\infty
\end{align*}\footnote{Inspired by the fact that for \(\int H_s\diff X_s = \int H_s\diff M_s+\int H_s\diff A_s\) and \(\mathbb{E}[\left(\int H_s\diff M_s+\int H_s\diff A_s\right)^2]\leq 2\mathbb{E}[\left(\int H_s \diff M_s\right)^2]+2\mathbb{E}[\left(\int H_s \diff A_s\right)^2]\leq2\cdot 4\mathbb{E}[\left(\int (H_s)^2 \diff [M]_s\right)^2]+2\mathbb{E}[\left(\int H_s \diff A_s\right)^2]\) from \(|a+b|^2\leq 2|a|^2+2|b|^2\) and BDG with \(C_2 = 4\).}\textcolor{blue}{One would still need to check that \(\lVert H\rVert_{L^2(M,A)}\) defines a norm}
where in our case \([X_0 + A + M]=[M]\). With this we have 
\begin{equation}
    \lVert\int_0^T\vartheta_s\diff X_s\rVert_{L^2}\leq K_X\lVert \vartheta \rVert_{L^2(M,A)}
\end{equation}
Indeed,
\begin{align*}\label{eq:bounding the L^2 norm of the stochastic integral}
    &\mathbb{E}[\left(\int_0^T H_s\diff X_s\right)^2]=\mathbb{E}[\left(\int_0^T H_s\diff M_s+\int_0^T H_s\diff A_s\right)^2]
    \\&
    \leq 2\mathbb{E}[\left(\int_0^T H_s \diff M_s\right)^2]+2\mathbb{E}[\left(\int_0^T H_s \diff A_s\right)^2]
    \\&
    \leq2\cdot 4\mathbb{E}[\int_0^T (H_s)^2 \diff [M]_s]+2\mathbb{E}[\left(\int_0^T H_s \diff A_s\right)^2]\leq
    \\&
     \leq2\cdot 4\mathbb{E}[\int_0^T (H_s)^2 \diff [M]_s]+2\mathbb{E}[\left(\int_0^T |H|_s \diff |A|_s\right)^2]\leq
    \lVert H\rVert_{L^2(M,A)}
\end{align*}
from \(|a+b|^2\leq 2|a|^2+2|b|^2\) and BDG with \(C_2 = 4\)\footnote{\(\mathbb{E}[\left(\int_0^TH_s\diff M_s\right)^2]\leq\mathbb{E}[\sup_{0\leq t\leq T}\left(\int_0^t H_s\diff M_s\right)^2]\leq C_2\mathbb{E}[\int_0^T(H_s)^2\diff [M]_s]\)}. Then we have
\begin{align}
    &\vartheta_n\overset{L^2(M,A)}{\to}\vartheta\implies\int_0^T\vartheta_s^n\diff X_s \overset{L^2}{\to}\int_0^T\vartheta_s\diff X_s
\end{align}
In this approach, we need to show that \(\lVert\cdot \rVert_{L^2(M,A)}\) is indeed a norm (or a semi-norm), that \(\mathcal{S}(\psi)\) is dense in \(L^2(M,A)\), and whether \(L^2(M,A)\) is rich enough to be interesting.
\section{Proof Attempt}
We present the first version of a generalization of \parencite[Thm~5.35]{Arandjelovic2025}. From \parencite[Thm.~5.32.]{Arandjelovic2025} we know that \(\mathcal{E}(\psi)\) is dense in \(L^p(M)\) for \(p\in[1,\infty)\) and \(M\in\mathcal{H}_0^{p+\epsilon}\) for some \(\epsilon>0\). 
Now consider a semimartingale \(X=X_0+M+A\), where \(M\in\mathcal{H}_0^{p+\epsilon}\) and \(A\) are adapted processes with RCLL trajectories, \(A_.(\omega)\) of finite variation and \(A_0=0\). Additionally, assume the following:
\begin{itemize}
    \item \(A<<\langle M\rangle\) path-wise (i.e. for every \(\omega\) as Lebesgue Stieltjes measures, with density process \(\lambda\).
    \item \(\operatorname*{ess\sup}\left(\langle M\rangle_{\infty}(\lambda^2+1)^*_{\infty}\right)<\infty\) where \((\lambda^2+1)^*_{\infty}=\sup_{t\geq 0}(\lambda^2+1)\).\footnote{Replace \(\infty\) by \(T\) in the finite horizon setting.}
\end{itemize}
Then we have, by density, of \(\mathcal{E}(\psi)\) in \(L^2(M)\) that
\begin{equation}
    \forall H\in L^2(M), \exists (H_n)_{n\in\mathbb{N}}\subset\mathcal{E}(\psi):\lim_{n\to\infty}\lVert H-H_n\rVert_{L^2(M)}=0\;.
\end{equation}
We can bound the \(L^2\) difference of the integrals by 
\begin{align}
    &\mathbb{E}[(H\bullet X)_{\infty}^2]\leq 2\left((\mathbb{E}[(H\bullet \langle M\rangle)_{\infty}^2]+\mathbb{E}[(H\bullet A)_{\infty}^2]\right)\\
    &\leq 2\left((\mathbb{E}[\langle M\rangle_{\infty}(H^2\bullet \langle M\rangle)_{\infty}]+\mathbb{E}[(H\lambda\bullet \langle M\rangle)_{\infty}^2]\right)\\
    &\leq2\left((\mathbb{E}[\langle M\rangle_{\infty}(H^2\bullet \langle M\rangle)_{\infty}]+\mathbb{E}[\langle M\rangle_{\infty}(H^2\lambda^2\bullet \langle M\rangle)_{\infty}]\right)\\
    &=2\mathbb{E}[\langle M\rangle_{\infty}(H^2(1+\lambda^2)\bullet\langle M\rangle)_{\infty}]\\
    &\leq 2\mathbb{E}[\langle M\rangle_{\infty}(1+\lambda^2)^{*}_{\infty}(H^2\bullet\langle M\rangle)_{\infty}]\\
    &\leq2\operatorname*{ess\sup}(\langle M\rangle_{\infty}(1+\lambda^2)^{*}_{\infty})\cdot\mathbb{E}[H^2\bullet \langle M\rangle]\\
    &=C(X)\lVert H\rVert_{L^2(M)}^2\;.
\end{align}
We used Cauchy Schwarz in step 2 and 3, along with the absolute continuity of \(A<<\langle M\rangle\). In step 5, we used the Hölder inequality with \(p=1,q=\infty\).\\
By assumption \(C(X)<\infty\), hence substituting for \(H\), \(Z^n=H-H^n\), we get that 
\begin{equation}
    \lim_{n\to\infty}\mathbb{E}[(Z^n\bullet X)_{\infty}^2]\leq\lim_{n\to\infty}C(X)\lVert H\rVert_{L^2(M)}^2=0
\end{equation}
By the argument given earlier, this allows for a deep mean-variance hedging result for a subclass of semimartingales.
\section{Different Approach}
We present a slightly different approach based on the observation \cref{eq:bounding the L^2 norm of the stochastic integral}. Notice that under the assumption that \(A<<\langle M\rangle\) path-wise (i.e. for every \(\omega\) as Lebesgue Stieltjes measures, with density process \(\lambda\) we have:
\begin{align*}
&\lVert H\rVert_{L^2(M,A)}^2=\mathbb{E}[\int_0^\infty (H_s)^2\diff [M]_s]+\mathbb{E}[\left(\int_0^\infty |H_s||\lambda_s|\diff[M]_s\right)^2]\\&\leq\mathbb{E}[\int_0^\infty (H_s)^2\diff [M]_s+\mathbb{E}[\int_0^{\infty}(H_s)^2\diff [M]_s\cdot\int_0^{\infty}(\lambda_s)^2\diff [M]_s]
\\&
\leq\mathbb{E}[\int_0^\infty (H_s)^2\diff [M]_s+\mathbb{E}[\left(\int_0^{\infty}(H_s)^2\diff [M]_s\right)^{p}]^{1/p}\cdot\mathbb{E}[\left(\int_0^{\infty}(\lambda_s)^2\diff [M]_s\right)^{q}]^{1/q}
\;,
\end{align*}
where we used the fact that if \(A<<[M]\) with density \(\lambda\) then \(|A|<<[M]\) with density \(|\lambda|\). Furthermore we used the Kunita-Watanabe inequality in the second step and Hölder's inequality in the third step. In the case \(p=1,q=\infty\) the above reads
\begin{equation}
    \lVert H\rVert_{L^2(M,A)}\leq\lVert H\rVert_{L^2(M)}\left(1+\operatorname*{ess\sup}\left(\int_0^{\infty}(\lambda_s)^2\diff [M]_s\right)\right)\;.
\end{equation}
Notice \(\int_0^{\infty}(\lambda_s)^2\diff [M]_s=\int_0^{\infty}|\lambda_s|\diff A_s\leq \int_0^{\infty}|\lambda_s|\diff |A|_s\). We impose the additional assumption
\begin{equation}
    \operatorname*{ess\sup}\int_0^{\infty}|\lambda_s|\diff |A|_s<\infty\;.
\end{equation}
Then, because we have the trivial inequality
\begin{align}\label{eq:essential boundedness assumption}
    \mathbb{E}[\int_0^{\infty}H_s^2\diff[M]_s]\leq  \mathbb{E}[\int_0^{\infty}H_s^2\diff[M]_s]+  \mathbb{E}[\left(\int_0^{\infty}|H_s||\lambda_s|\diff[M]_s\right)^2]=\lVert H\rVert_{L^2(M,A)}
\end{align}
we have that the norms on \(L^2(M)\) and \(L^2(M,A)\),under the condition that \(X = X_0+M+A\) with \(A<<[M]\) and \(M\in\mathcal{H}^{p+\epsilon}_0\) for some \(p\in[0,\infty)\) and \(\epsilon>0\), along with the additional assumption \cref{eq:essential boundedness assumption}, are equivalent. Because both spaces \(L^2(M)\) and \(L^2(M,A)\) are subspaces of the space of \(\mathcal
P\)-measurable functions in, where \(\mathcal{P}\) is the predictable \(\sigma\)-algebra on \(\Omega\times[0,\infty)\), they are defined through their norms; the spaces are equal:
\begin{equation}
    L^2(M) =L^2(M,A)\;.
\end{equation}
Hence, density of \(\mathcal{E}(\psi)\) in \(L^2(M)\) is obviously equivalent to density of \(\mathcal{E}(\psi)\) in \(L^2(M,A)\) under the stated assumptions on \(M\) and \(X\). Hence, for every \(H\in L^2(M,A)\) there exists \((H_n)_{n\in\mathbb{N}}\subset \mathcal{E}(\psi)\) such that \(\lim_{n\to\infty}\lVert H-H_n\rVert_{L^2(M,A)}=0\). Then we have, by the observation \cref{eq:bounding the L^2 norm of the stochastic integral}, with \(T=\infty\) and \(H_s\coloneqq (H-H_n)_s\)
\begin{equation}
    \lVert (H-H^n)\bullet X\rVert_{L^2}
    \lesssim\lVert H-H^n\rVert_{L^2(M,A)}\;,
\end{equation}
thus
\begin{equation}
    \lim_{n\to\infty}\lVert (H-H^n)\bullet X\rVert_{L^2}=0
\end{equation}
\bigbreak
The two approaches differ in their additional assumption
\begin{enumerate}
    \item \(\operatorname*{ess\sup}(\langle M\rangle_{\infty}(1+\lambda^2)^{*}_{\infty})<\infty\)
    \item  \(\operatorname*{ess\sup}\int_0^{\infty}|\lambda_s|\diff |A|_s<\infty\)
\end{enumerate}
Notice that
\begin{equation}
    \int_0^\infty(\lambda_s)^2\diff [M]_s\leq\int_0^{\infty}(1+\lambda_s^2)\diff[M]_s\leq(1+\lambda_s^2)_{\infty}^*[M]_{\infty}\,
\end{equation}
where we used the non-negativity and monotonicity of the process \([M]\). Hence, (2.) is a weaker assumption than (1.).
\textcolor{blue}{potential problem with this approach. The norms are equivalent but not really norms. First one needs to check that \(\lVert \cdot\rVert_{L^2(M,A)}\) defines at least a subadditiv, homogenous map and then that the process which are identified by \(\lVert X-Y\rVert_{L^2(M,A)}=0\) are also identified in \(L^2(M)\), that is \(\lVert X-Y\rVert_{L^2(M)}=0\). So that the norms are norms on the same quotient space under some equivalence relation identifying processes}
First, notice that \(\lVert H\rVert_{L^2(M,A)}=0\) is subadditive
Under the assumption that \(A<<[M]\) has density \(\lambda\), we have that \(\lVert X-Y\rVert_{L^2(M)}=0\implies\lVert X-Y\rVert_{L^2(M,A)}=0\). This follows from the fact that \(c\lVert X\rVert_{L^2(M)}\leq\lVert X\rVert_{L^2(M,A)}\leq C\lVert X\rVert_{L^2(M)}\) for \(c=1,C=(1+\operatorname*{ess\sup}\int_0^{\infty}|\lambda_s|\diff |A|_s)\). By assumption, \(c,C\) are finite. This implies that under the equivalence relation \(X\sim Y\) for \(X\sim Y\) for \(\lVert X-Y\rVert_{L^2(M,A)}=0\). Hence, they are norms on the same space \(L^2(\Omega\times [0,T],\mathcal{P})/\sim\), and in particular, it makes sense to say that they are equivalent. 
\bigbreak
This proves the following result:
\begin{theorem}[Semimartingale Deep Hedging]\label{thm:general deep hedging}
    Let \(X=X_0+M+A\) be a semimartingale with \(M\in \mathcal{H}_{0}^{2+\epsilon}\) for some \(\epsilon>0\), \(A<<[M]\) with density \(\lambda=(\lambda_s)_{s\in\mathbb{R}_+}\) and with \(\operatorname*{ess\sup}\int_0^{\infty}|\lambda_s|\diff |A|_s<\infty\) \. Let \(H\in L^2(\mathbb{P})\) and \((V_0^*,\vartheta^*)\) be the solution to \cref{thm:Hilbert Projection for Mean Sqaure Hedging}. Then there exists \((\vartheta_n)_{n\in\mathbb{N}}\subset\mathcal{E}(\mathcal{\psi})\) such that:
    \begin{equation}
        \min_{\vartheta\in \Theta}\lVert H-V_0^*-\int_{0}^T\vartheta_s\diff X_s\rVert_{L^2}=\inf_{\vartheta\in\mathcal{E}(\psi)}\lVert H-V_0^*-\int_{0}^T\vartheta_s\diff X_s\rVert_{L^2}\;,
    \end{equation}
\end{theorem}
where \(\lVert H-V_0^*-\int_{0}^T\vartheta^*_s\diff X_s\rVert_{L^2}=\min_{\vartheta\in \Theta}\lVert H-V_0^*-\int_{0}^T\vartheta_s\diff X_s\rVert_{L^2}\).
\begin{proof}
    The proof is given above and relies on the fact that in this case \(L^2(M)=L^2(M,A)\). By \parencite[Thm.~5.32]{Arandjelovic2025} \(\mathcal{E}(\psi)\) is dense in \(L^2(M)\) and hence also trivially in \(L^2(M,A)\) as well. The only difference is that the optimal \(V_0^*,H^*\) are for the more general problem with \(X= X_0+M+A\). The proof is the same as in \parencite[THm.~5.35]{Arandjelovic2025} as for any \(\delta>0\) there exists \(U\in\mathcal{E}(\psi)\) such that \(\lVert(U^*\bullet X)_{\infty}-(U\bullet X)_{\infty}\rVert_{L^2}\leq \lVert U^*-U\rVert_{L^2(M,A)}<\delta\) by the calculation in \cref{eq:bounding the L^2 norm of the stochastic integral}.
\end{proof}
One should give examples, if they exist, of finite variation processes \(A\) satisfying the assmuptions:
\begin{itemize}
    \item \(A<<[M]\) with density \(\lambda\)
    \item \(\operatorname*{ess\sup}\int_0^{\infty}|\lambda_s|\diff |A|_s<\infty\)
\end{itemize}

Furthermore we will think of a more general approach generalizing \parencite[Thm.~5.35]{Arandjelovic2025} again with \(X\) a semimartingale and approximants from \(\mathcal{S}(\psi)\). For this we might think again about strategy (1) and if there are sensible conditions on \(\vartheta\) that transfer to conditions on \(\vartheta_n\) such that the respective family of stochastic integrals (more precisely the square processes) become a uniformly integrable family.

\textcolor{blue}{Possible, non-trivial process satisfying assumptions: Let}
\begin{equation}
    S_t = \lambda B_t + \mu t
\end{equation}
Then \(\langle S\rangle_t= \lambda^2 t\) and \(A_t=\mu t\).
\end{document}
% so it automatically matches the exact same layout/preamble.

\section{Proof attempts}

% Put your drafts here. Example skeleton:
%
%\begin{proof}[Proof attempt (Lemma~\ref{lem:...})]
%    ...
%\end{proof}
\section{Extension to semi-martingale case}
\cref{thm:cont-deep-ms-hedging-martingale} was in the realm of classical martingale stochastic calculus, with integrators \(M\in\mathcal{H}^{2}\)(even \(\mathcal{H}^{2+\epsilon}\) and integrands \(\vartheta\in L^2(M)\). In light of \cref{thm:Approx of stoch. Integrals}, we aim to prove a similar result as in \cref{thm:cont-deep-ms-hedging-martingale} for the more general case where \(M\) is a semimartingale. Following \parencite{Schweizer2001GuidedTourQuadraticHedging} we restrict the class of possible strategies to \(\vartheta_{GLP}\) such that \(G_T(\vartheta_{GLP})\) is a closed subspace in \(L^2(\mathbb{P})\). In this case, we know that the minimization problem has a solution, as \(\mathcal{A}=\mathbb{R}+G_T(\Theta_{GLP})\) is a closed subspace in \(L^2(\mathbb{P})\).; this means there exists an optimal trading strategy \((V_0^*,\vartheta^*)\in\mathcal{A}=\mathbb{R}+G_T(\Theta_{GLP})\) (\textcolor{blue}{I guess we don't even need strong conditions to get \(\mathcal{A}\) closed as we don't need explicit expression of the optimal trading strategy at first, hence we can directly project on \(\bar{\mathcal{A}}\) which is closed in \(L^2(\mathbb{P})\). In that line of thought, any \(\Theta\subset\Theta_2\) is sufficient for the existence of a solution in \(\bar{\mathcal{A}}\)}.) Then we have the following consequence of the Hilbert Projection theorem
\begin{theorem}[Hilbert Projection for Mean Sqaure Hedging]\label{thm:Hilbert Projection for Mean Sqaure Hedging}
    Let \(\Theta\subset\Theta_2\) and \(\mathcal{A}=\mathbb{R}+G_T(\Theta)\subset L^2(\mathbb{P})\). Then the \(L^2\)-closure \(\bar{\mathcal{A}}\) is a closed subspace of \(L^2(\mathbb{P})\). Hence, there exists \((V_0^*,\vartheta^*)\) such that 
    \begin{equation}
   \lVert H-V_0^*-\int_{0}^T\vartheta^*_s\diff X_s\rVert_{L^2}= \inf_{(c,\vartheta)\in\bar{\mathcal{A}}}\lVert H-c-\int_{0}^T\vartheta_s\diff X_s\rVert_{L^2}
    \end{equation}
    and the solution is unique. 
\end{theorem}
\begin{remark}
    Notice, we can study the related problem where we fix \(V_0\) (corresponds to the situation where the price \(V_0\) of the liability \(H\) might be determined by the market) and study
    \begin{equation}
        \inf_{\vartheta\in\Theta}\lVert H-V_0-\int_{0}^T\vartheta_s\diff X_s\rVert_{L^2}
    \end{equation}
    Again we can consider the the \(L^2\) closure of \(G_T(\Theta)=\left\{\int_0^T \vartheta_s\diff X_s:\vartheta\in\Theta\right\}\) and use the Hilbert Projection theorem on this closed subspace of \(L^2\).
\end{remark}
Then, crucially, notice that from \cref{thm:Approx of stoch. Integrals} we obtain the convergence in probability of the final values of the respective integral processes. Indeed, the \(\mathbb{P}\)-\(\mathbb{P}\)-ucp convergence of the integral processes, restricted to a finite time horizon \(T\), reads as:
\begin{equation}
    \sup_{0\leq t\leq T}\left|\int_{0}^t\vartheta_s^n\diff X_s-\int_{0}^t\vartheta^*_s\diff X_s\right|\to 0\quad\text{in Probability}
\end{equation}
where \(\{\vartheta^n\}_{n\in\mathbb{N}}\in\mathcal{S}(\psi)\) is a minimizing sequence from \cref{thm:Approx of stoch. Integrals}
\\\\
In particular, this implies 
\begin{equation}
    \int_{0}^T\vartheta_s^n\diff X_s\xrightarrow{\mathbb{P}}\int_{0}^T\vartheta^*_s\diff X_s\;,
\end{equation}
where convergence holds in probability. Imposing the additional assumption that \(\left\{\left|\int_{0}^T\vartheta_s^n\diff X_s\right|^2\right\}_{n\in\mathbb{N}}\) is a uniformly integrable family of random variables  leads to the \(L^2\) convergence\footnote{In particular, the condition also implies that the random variable \(\int_{0}^T\vartheta_s^n\diff X_s\in L^2(\mathbb{P)}\quad\forall n\in\mathbb{N}\)}: 
\begin{equation}
\int_{0}^T\vartheta_s^n\diff X_s \xrightarrow{L^2(\mathbb{P})} \int_{0}^T\vartheta^*_s\diff X_s\;.
\end{equation}
We do NOT want the uniform integrability of \(\left\{\left|\int_{0}^T\vartheta_s^n\diff X_s\right|^2\right\}_{n\in\mathbb{N}}\) to follow from conditions on a sequence of trading strategies. Indeed, \cref{thm:Approx of stoch. Integrals} asserts the existence of a sequence \(\{\vartheta^n\}_{n\in\mathbb{N}}\subset\mathcal{S}(\psi)\). We don't want a theorem that says something like: 
\\\\
"If that sequence happens to be square uniformly integrable, then it follows...."
\\\\
Hence, if necessary, we shall impose conditions only on the integrator \(X\). Recall that for \(\vartheta_n\in\mathcal{S}(\psi)\) that is \(\vartheta^n_t=\varphi_0\mathbbm{1}_{\{0\}}+\sum_{i=1}^l\varphi_i^n\mathbbm{1}_{\tau_i,\tau_{i+1}}\) and \(\varphi_i^{(n)}=f^{(n)}(Y^{\tau_i}_{t_1},\dots,Y^{\tau_i}_{t_r})\) where \(f\in\mathcal{NN}((\mathbb{R}^{d\times r})^*,\psi)\)\footnote{the parameter \(r\in\mathbb{N}\) is not further specified. It relies to the existence result of \parencite[Thm.~5.23.]{Arandjelovic2025}}, the stochastic integral takes the form \(\int_0^T\vartheta^n_s\diff X_s = \sum_{i=1}^k \varphi^{(n)}_i(X_{\tau_i\wedge T}-X_{\tau_{i+1}\wedge T})\). Note that \(\left|\int_0^T\vartheta^n_s\diff X_s\right|^2 \leq \sum_{i=1}^k \left|\varphi^{(n)}_i\right|^2\left|(X_{\tau_i\wedge T}-X_{\tau_{i+1}\wedge T})\right|^2\). But there can't exist a uniform bound (in n) on \(|\varphi_i^n|\) as this would imply that neural networks can't approximate constants above this bound which is absurd. Also, notice for \(\tau<\sigma\) two stopping times, that \(f_n = c_n\mathbbm{1}_{[\tau,\sigma)}\in\mathcal{S}(\psi)\) for all \(c_n\in\mathbb{R}\). Take \(c_n\to\infty\) as \(n\to\infty\) a real deterministic sequence, then \(\int_0^T(f_n)_s\diff X_s=c_n(X_{\tau\wedge T}-X_{\sigma\wedge T})\), and hence \(\mathbb{E}[\left|\int_0^T(f_n)_s\diff X_s\right|^2]=\mathbb{E}[\left|c_n\right|^2\left|(X_{\tau\wedge T}-X_{\sigma\wedge T})\right|^2]\). If this is to be UI, then in particular, it is bounded in \(L^1\). But \(\sup_{n\in\mathbb{N}}\mathbb{E}[c_n^2(X_{\tau\wedge T}-X_{\sigma\wedge T})^2]=(\sup_{n\in\mathbb{N}}c_n^2)\mathbb{E}[(X_{\tau\wedge T}-X_{\sigma\wedge T})^2]<\infty\) by assumption. But \(\sup_{n\in\mathbb{N}}c_n^2=\infty\) hence we must have \(\mathbb{E}[(X_{\tau\wedge T}-X_{\sigma\wedge T})^2]=0\) which implies \(X_{\tau\wedge T}-X_{\sigma\wedge T}=0\quad\mathbb{P}\)-a.s. As this holds for all \(\mathbb{F}\)-stopping times \(\tau,\sigma\), in particular deterministic ones, we have \(X\equiv 0\) almost surely. Hence, there can't be a (non-trivial) condition on \(X\) making the family \(\left\{\left|\int_{0}^T\vartheta_s^n\diff X_s\right|^2\right\}_{n\in\mathbb{N}}\) UI, as the latter being UI for any sequence \((\vartheta^n)_{n\in\mathbb{N}}\) implies that \(X\equiv 0\).
\bigbreak
\textcolor{blue}{As an initial observation, recall that \(\int_{0}^T\vartheta_s^n\diff X_s\xrightarrow{\mathbb{P}}\int_{0}^T\vartheta^*_s\diff X_s\;\) implies that there exists a subsequence \(\left\{\int_{0}^T\vartheta_s^{n(k)}\diff X_s\right\}_{k\in\mathbb{N}}\) that converges \(\mathbb{P}\)-a.s. to \(\int_{0}^T\vartheta_s^{*}\diff X_s\). Hence, on a full measure set, it is (uniformly) bounded as a real converging sequence (assuming that \(\int_{0}^T\vartheta_s^{*}\diff X_s<\infty\) on that full measure set). If it is only \(\mathbb{P}\)-a.s. finite, take the intersection of the respective full measure sets to reach the same conclusion.) Hence, there exists a positive constant \(M\in\mathbb{R}_+\) such that \(\left|\int_{0}^T\vartheta_s^{n(k)}\diff X_s\right|^2<M^2\quad \mathbb{P}\)-a.s. for all \(k\in\mathbb{N}\). But \(M^2\) is integrable; hence, the family \(\left\{\left|\int_{0}^T\vartheta_s^{n(k)}\diff X_s\right|^2\right\}_{k\in\mathbb{N}}\) is uniformly integrable.\footnote{If \(H\geq |X_n|\) for all \(n\in\mathbb{N}\), we have \(\mathbb{E}[|X|\mathbbm{1}_{|X|\geq K}]\leq\mathbb{E}[H\mathbbm{1}_{H\geq K}]\) for any \(K>0\). As the left hand side going to 0 as \(K\to\infty\) it gives us the uniform integrability of \(\{X_n\}_{n\in\mathbb{N}}\)}. Thus, we get }
\\\\
\textcolor{red}{This conclusion is WRONG. Notice that the bound \(M = M(\omega)\) and hence doesn't define a random variable \(H\in L^2\)(at least not because it is constant almost surely). Indeed the argument used, goes as follows: For every \(\omega\) the sequence \(\left(\int_{0}^T\vartheta_s^{n(k)}\diff X_s\right)(\omega)\) is bounded as a real converging sequence, hence \(\left|\left(\int_{0}^T\vartheta_s^{n(k)}\diff X_s\right)(\omega)\right|\leq M(\omega)\).}
\\\\
\begin{equation}
\int_{0}^T\vartheta_s^{n(k)}\diff X_s \xrightarrow{L^2(\mathbb{P})} \int_{0}^T\vartheta^*_s\diff X_s\;\text{ as } k\to\infty
\end{equation}
Hence, restricting ourselves to this subsequence (which obviously still converges in probability to the same limit), we have 
\begin{align*}
    &\inf_{k\in\mathbb{N}}\Vert[H-V_0^*-\int_{0}^T\vartheta_s^{n(k)}\diff X_s]\rVert_{L^2}\leq \liminf_{k\to\infty}\lVert H-V_0^*-\int_{0}^T\vartheta_s^{n(k)}\diff X_s]\rVert_{L^2}\\&\leq \liminf_{k\to\infty}\lVert H-V_0^*-\int_{0}^T\vartheta_s^{*}\diff X_s\rVert_{L^2}+\lVert \int_{0}^T\vartheta_s^{*}\diff X_s -\int_{0}^T\vartheta_s^{n(k)}\diff X_s\rVert_{L^2}
    \\&=\lVert H-V_0^*-\int_{0}^T\vartheta_s^{*}\diff X_s\rVert_{L^2}
\end{align*}
where \(\lVert H-V_0^*-\int_{0}^T\vartheta_s^{*}\diff X_s\rVert_{L^2}=\inf_{\vartheta\in\Theta}\lVert H-V_0^*-\int_{0}^T\vartheta_s\diff X_s\rVert_{L^2}\).
By \cref{rem:inf and sequences}, the existence of this sequence is equivalent to a more general statement:
\begin{theorem}[Semimartingale Deep Hedging]\label{thm:general deep hedging}
    Let \(X\) be a semimartingale and \(\Theta\subset\Theta_2\). Let \((V_0^*,\vartheta^*)\) be the solution to \cref{thm:Hilbert Projection for Mean Sqaure Hedging}. Then there exists \((\vartheta_n)_{n\in\mathbb{N}}\subset\mathcal{S}(\mathcal{\psi})\) such that:
    \begin{equation}
        \min_{\vartheta\in \Theta}\lVert H-V_0^*-\int_{0}^T\vartheta_s\diff X_s\rVert_{L^2}=\inf_{\vartheta\in\mathcal{S}(\psi)}\lVert H-V_0^*-\int_{0}^T\vartheta_s\diff X_s\rVert_{L^2}\;,
    \end{equation}
\end{theorem}
where \(\lVert H-V_0^*-\int_{0}^T\vartheta^*_s\diff X_s\rVert_{L^2}=\min_{\vartheta\in \Theta}\lVert H-V_0^*-\int_{0}^T\vartheta_s\diff X_s\rVert_{L^2}\).
\\\\
\newpage
\underline{Possible strategies}:
In order to get deep-hedging result similar to \cref{thm:general deep hedging} we need 
\begin{equation}
    \int_0^T\vartheta^n_s\diff X_s \overset{L^2}{\to} \int_0^T\vartheta_s\diff X_s
\end{equation}
For an approximating sequence \((\vartheta_n)_{n\in\mathbb{N}}\subset\mathcal{S}(\psi)\) of \(\vartheta\in\tilde{\Theta}\), where the sense of approximation is to be determinded and the set \(\tilde{\Theta}\) such that \(\mathcal{S}(\psi)\subset\tilde{\Theta}\). Assuming this is done, we have
\begin{align*}
  &\inf_{\vartheta\in\mathcal{S}(\psi)}\lVert H-V_0^*-\int_0^T \vartheta_s^n\diff X_s\rVert_{L^2}\\&\leq
    \inf_{n\in\mathbb{N}}\lVert H-V_0^*-\int_0^T \vartheta_s^n\diff X_s\rVert_{L^2}\leq \liminf_{n\to\infty}\lVert H-V_0^*-\int_0^T \vartheta_s^n\diff X_s\rVert_{L^2}
    \\&\leq \liminf_{n\to\infty}\lVert H-V_0^*-\int_0^T \vartheta_s^*\diff X_s\rVert_{L^2}+\lVert \int_0^T \vartheta_s^n\diff X_s-\int_0^T \vartheta_s^*\diff X_s\rVert_{L^2}\\&=\lVert H-V_0^*-\int_0^T \vartheta_s^*\diff X_s\rVert_{L^2}\;,
\end{align*}
where \(\lVert H-V_0^*-\int_0^T \vartheta_s^*\diff X_s\rVert_{L^2}=\argmin_{\vartheta\in\Theta}\lVert H-V_0^*-\int_0^T \vartheta_s^*\diff X_s\rVert_{L^2}\) for some set \(\Theta\subset\Theta_2\). On the other hand we have \(\inf_{\vartheta\in\tilde{\Theta}}\lVert H-V_0^*-\int_0^T \vartheta_s^n\diff X_s\rVert_{L^2}\leq\inf_{\vartheta\in\mathcal{S}(\psi)}\lVert H-V_0^*-\int_0^T \vartheta_s^n\diff X_s\rVert_{L^2}\), because \(\mathcal{S}(\psi)\subset\tilde{\Theta}\). Thus we get equality. \\
We consider two possible strategies.
\bigbreak
(1): From \parencite[Thm~5.29.]{Arandjelovic2025} we have \(\mathcal{S}(\psi)\) is dense in \(\mathcal{S}\) which is dense in \(\mathcal{L}\), the space of càglàd processes (see Protter for this). Then for \(\vartheta_n\overset{\mathbb{P}-ucp}{\to}\vartheta\) where \(\vartheta\in \mathcal{L}\), we have 
\begin{align*}
    \int \vartheta_s^n\diff X_s\overset{\mathbb{P}-ucp}{\to}\int \vartheta_s\diff X_s\implies\int_0^T \vartheta_s^n\diff X_s\overset{\mathbb{P}}{\to}\int_0^T\vartheta_s\diff X_s
\end{align*}
and 
\begin{align*}
    &\left(\vartheta_s^n\diff X_s\overset{\mathbb{P}}{\to}\int_0^T\vartheta_s\diff X_s \right) + \left(\left\{\left|\int_{0}^T\vartheta_s^{n}\diff X_s\right|^2\right\}_{n\in\mathbb{N}}\text{uniformly integrable}\right)\\&\implies \left(\int_0^T\vartheta_s^n\diff X_s\overset{L^2}{\to}\int_0^T\vartheta_s\diff X_s \right)
\end{align*}
WARNING: Notice that for sequences of the form \(\vartheta_n=c_n\cdot\mathbbm{1}_{[\tau,\sigma)}\in\mathcal{S}(\psi)\) for \(\lim_{n\to\infty}c_n=\infty\), requiring that \(\left\{\left|\int_{0}^T\vartheta_s^{n}\diff X_s\right|^2\right\}_{n\in\mathbb{N}}\) is uniformly integrable implies \(X=0\) as a process. This is a pathological example as it doesn't converge to an interesting limiting process \(\vartheta\), which we assume. This just shows that  \(\left\{\left|\int_{0}^T\vartheta_s^{n}\diff X_s\right|^2\right\}_{n\in\mathbb{N}}\) can't be UI for all \((\vartheta_n)_{n\in\mathbb{N}}\subset\mathcal{S}(\psi)\). 
\bigbreak
(2): Notice that \(Y_n\overset{L^2-ucp}{\to}Y\) i.e. \(\lim_{n\to\infty}\mathbb{E}[\sup_{0\leq t\leq T}|Y_t^n-Y_t|^2]=0\) implies \(Y^n_T\overset{L^2}{\to}Y_T\). Hence, we want a version of \parencite[Thm~5.29]{Arandjelovic2025} with \(L^2-ucp\) convergence of the integrals. For this, we might want to use the fact that for \(\Phi(\cdot)=(\cdot)^2\), \parencite[Thm~5.28]{Arandjelovic2025} implies that \(\mathcal{S}(\psi)\) is \(L^2-ucp\) dense in \(\mathcal{S}\) and use some kind of \(L^2-ucp\) continuity of the stochastic integral \(H\to H\bullet X\) for well chosen classes for \(H\) and \(X\). An clearly weaker assumption would be to ask for convergence in \(\mathcal{H}^2\) which implies the \(L^2\) convergence of the end-points (Doob). A well known case where we have the continuity of the integral map w.r.t. to the \(L^2\)-convergence of the end-points is the isometry property (Itô isomtrey) of the map 
\begin{align}
    \begin{cases}
        H\to H\bullet M\\
        L^2(M)\to\mathcal{H}^2_0
    \end{cases}
\end{align}
for \(M\in\mathcal{H}^2_{0,loc}\)
Indeed in this case, by \parencite[Thm~5.32]{Arandjelovic2025} we even have that \(\mathcal{E}(\psi)\) is dense in \(L^p(M)\) for \(p\in[1,\infty)\) if \(M\in\mathcal{H}_0^{p+\epsilon}\) for some \(\epsilon>0\). (Then we also have that \(\mathcal{S}(\psi)\) is dense in \(L^p(M)\)). Hence we have 
\begin{equation}
    \lVert H\bullet M\rVert_{\mathcal{H}^2}=\lVert H\rVert_{L^2(M)}
\end{equation}
Using the sequence \(H_n=\vartheta_n-\vartheta\) and the fact that convergence in \(\mathcal{H}^2\) implies convergence in \(L^2-ucp\) (by Doob's maximal inequality, which is applicable because \(H\bullet M\) is a square integrable martingale in this case) we recover \parencite[Thm~5.35]{Arandjelovic2025}.
\\\\
To get a slight extension of \parencite[Thm~5.35]{Arandjelovic2025}, we notice that we don't need the full isometry property; rather, we need the boundedness (continuity) of the map \(H\to H \bullet M\) and the corresponding density of \(\mathcal{S}(\psi)\). Hence, for a semimartingale 
\begin{equation}
    X = X_0 + A + M
\end{equation}
for \(A\) a finite-variation process and \(M\in\mathcal{H}^2_{0,loc}\) (the addition of \(A\) destroys the Itô-isometry property) we try to bound 
\begin{equation}
     \lVert \left(H\bullet X\right)_T\rVert_{L^2}\leq K_X\lVert H\rVert_{L^2(M)}
\end{equation}
for some constant \(K_X = K(X)\). By \(|a+b|^2\leq 2|a|^2+2|b|^2\) it is enough to bound \(\lVert H\bullet A\rVert_{\mathcal{H}^2}\) and \(\lVert H\bullet M\rVert_{\mathcal{H}^2}\) both by \(\lVert H\rVert_{L^2(M)}\) and some multiplicative constant. Notice that for \(M\in\mathcal{H}^2_{0,loc}\) and \( H \in L^2(M)\), \(H\bullet M\) is a square integrable martingale; hence, we can bound \(\lVert H\bullet M\rVert_{\mathcal{H}^2}=\lVert (H\bullet M)_T\rVert_{L^2}\leq 4\lVert H\rVert_{L^2}\) by Doob. For the finite variation part, we might need assumptions on \(A\) (and hence on \(X\)).
\\
Based on this oberservations we consider the following approach. Consider a new space \(L^2(M,A)\) for a semimartingale \(X=X_0+M+A\) where \(A\) is a finite-variation process and \(M\in\mathcal{H}^2_{0,loc}\). Define it by \footnote{Note that in our case \([X]=[M]\) and \(L^2(M,A)\subset L^2(M)\), so \(\int_0^T\vartheta_s\diff M_s\in\mathcal{H}_0^2\); hence, \(\left(\int \vartheta_s\diff M_s\right)^2-[\int \vartheta_s\diff M_s]\) is a martingale starting at 0, and \([\int\vartheta_s\diff M_s]=\int \vartheta_s^2\diff [M]_s\). Hence, \(\mathbb{E}[\left(\int_0^T \vartheta_s\diff M_s\right)^2]=\mathbb{E}[\int_0^T \vartheta_s^2\diff [M]_s]\)}
\begin{align*}
H\in L^2(M,A) \iff \left(\mathbb{E}[\int_0^T H_s^2\diff [X]_s]+\mathbb{E}[\left(\int_0^T |H|_s\diff |A|_s\right)^2]\right)^{\frac{1}{2}}\coloneqq\lVert H\rVert_{L^2(M,A)}<\infty
\end{align*}\footnote{Inspired by the fact that for \(\int H_s\diff X_s = \int H_s\diff M_s+\int H_s\diff A_s\) and \(\mathbb{E}[\left(\int H_s\diff M_s+\int H_s\diff A_s\right)^2]\leq 2\mathbb{E}[\left(\int H_s \diff M_s\right)^2]+2\mathbb{E}[\left(\int H_s \diff A_s\right)^2]\leq2\cdot 4\mathbb{E}[\left(\int (H_s)^2 \diff [M]_s\right)^2]+2\mathbb{E}[\left(\int H_s \diff A_s\right)^2]\) from \(|a+b|^2\leq 2|a|^2+2|b|^2\) and BDG with \(C_2 = 4\).}\textcolor{blue}{One would still need to check that \(\lVert H\rVert_{L^2(M,A)}\) defines a norm}
where in our case \([X_0 + A + M]=[M]\). With this we have 
\begin{equation}
    \lVert\int_0^T\vartheta_s\diff X_s\rVert_{L^2}\leq K_X\lVert \vartheta \rVert_{L^2(M,A)}
\end{equation}
Indeed,
\begin{align*}\label{eq:bounding the L^2 norm of the stochastic integral}
    &\mathbb{E}[\left(\int_0^T H_s\diff X_s\right)^2]=\mathbb{E}[\left(\int_0^T H_s\diff M_s+\int_0^T H_s\diff A_s\right)^2]
    \\&
    \leq 2\mathbb{E}[\left(\int_0^T H_s \diff M_s\right)^2]+2\mathbb{E}[\left(\int_0^T H_s \diff A_s\right)^2]
    \\&
    \leq2\cdot 4\mathbb{E}[\int_0^T (H_s)^2 \diff [M]_s]+2\mathbb{E}[\left(\int_0^T H_s \diff A_s\right)^2]\leq
    \\&
     \leq2\cdot 4\mathbb{E}[\int_0^T (H_s)^2 \diff [M]_s]+2\mathbb{E}[\left(\int_0^T |H|_s \diff |A|_s\right)^2]\leq
    \lVert H\rVert_{L^2(M,A)}
\end{align*}
from \(|a+b|^2\leq 2|a|^2+2|b|^2\) and BDG with \(C_2 = 4\)\footnote{\(\mathbb{E}[\left(\int_0^TH_s\diff M_s\right)^2]\leq\mathbb{E}[\sup_{0\leq t\leq T}\left(\int_0^t H_s\diff M_s\right)^2]\leq C_2\mathbb{E}[\int_0^T(H_s)^2\diff [M]_s]\)}. Then we have
\begin{align}
    &\vartheta_n\overset{L^2(M,A)}{\to}\vartheta\implies\int_0^T\vartheta_s^n\diff X_s \overset{L^2}{\to}\int_0^T\vartheta_s\diff X_s
\end{align}
In this approach, we need to show that \(\lVert\cdot \rVert_{L^2(M,A)}\) is indeed a norm (or a semi-norm), that \(\mathcal{S}(\psi)\) is dense in \(L^2(M,A)\), and whether \(L^2(M,A)\) is rich enough to be interesting.
\section{Proof Attempt}
We present the first version of a generalization of \parencite[Thm~5.35]{Arandjelovic2025}. From \parencite[Thm.~5.32.]{Arandjelovic2025} we know that \(\mathcal{E}(\psi)\) is dense in \(L^p(M)\) for \(p\in[1,\infty)\) and \(M\in\mathcal{H}_0^{p+\epsilon}\) for some \(\epsilon>0\). 
Now consider a semimartingale \(X=X_0+M+A\), where \(M\in\mathcal{H}_0^{p+\epsilon}\) and \(A\) are adapted processes with RCLL trajectories, \(A_.(\omega)\) of finite variation and \(A_0=0\). Additionally, assume the following:
\begin{itemize}
    \item \(A<<\langle M\rangle\) path-wise (i.e. for every \(\omega\) as Lebesgue Stieltjes measures, with density process \(\lambda\).
    \item \(\operatorname*{ess\sup}\left(\langle M\rangle_{\infty}(\lambda^2+1)^*_{\infty}\right)<\infty\) where \((\lambda^2+1)^*_{\infty}=\sup_{t\geq 0}(\lambda^2+1)\).\footnote{Replace \(\infty\) by \(T\) in the finite horizon setting.}
\end{itemize}
Then we have, by density, of \(\mathcal{E}(\psi)\) in \(L^2(M)\) that
\begin{equation}
    \forall H\in L^2(M), \exists (H_n)_{n\in\mathbb{N}}\subset\mathcal{E}(\psi):\lim_{n\to\infty}\lVert H-H_n\rVert_{L^2(M)}=0\;.
\end{equation}
We can bound the \(L^2\) difference of the integrals by 
\begin{align}
    &\mathbb{E}[(H\bullet X)_{\infty}^2]\leq 2\left((\mathbb{E}[(H\bullet \langle M\rangle)_{\infty}^2]+\mathbb{E}[(H\bullet A)_{\infty}^2]\right)\\
    &\leq 2\left((\mathbb{E}[\langle M\rangle_{\infty}(H^2\bullet \langle M\rangle)_{\infty}]+\mathbb{E}[(H\lambda\bullet \langle M\rangle)_{\infty}^2]\right)\\
    &\leq2\left((\mathbb{E}[\langle M\rangle_{\infty}(H^2\bullet \langle M\rangle)_{\infty}]+\mathbb{E}[\langle M\rangle_{\infty}(H^2\lambda^2\bullet \langle M\rangle)_{\infty}]\right)\\
    &=2\mathbb{E}[\langle M\rangle_{\infty}(H^2(1+\lambda^2)\bullet\langle M\rangle)_{\infty}]\\
    &\leq 2\mathbb{E}[\langle M\rangle_{\infty}(1+\lambda^2)^{*}_{\infty}(H^2\bullet\langle M\rangle)_{\infty}]\\
    &\leq2\operatorname*{ess\sup}(\langle M\rangle_{\infty}(1+\lambda^2)^{*}_{\infty})\cdot\mathbb{E}[H^2\bullet \langle M\rangle]\\
    &=C(X)\lVert H\rVert_{L^2(M)}^2\;.
\end{align}
We used Cauchy Schwarz in step 2 and 3, along with the absolute continuity of \(A<<\langle M\rangle\). In step 5, we used the Hölder inequality with \(p=1,q=\infty\).\\
By assumption \(C(X)<\infty\), hence substituting for \(H\), \(Z^n=H-H^n\), we get that 
\begin{equation}
    \lim_{n\to\infty}\mathbb{E}[(Z^n\bullet X)_{\infty}^2]\leq\lim_{n\to\infty}C(X)\lVert H\rVert_{L^2(M)}^2=0
\end{equation}
By the argument given earlier, this allows for a deep mean-variance hedging result for a subclass of semimartingales.
\section{Different Approach}
We present a slightly different approach based on the observation \cref{eq:bounding the L^2 norm of the stochastic integral}. Notice that under the assumption that \(A<<\langle M\rangle\) path-wise (i.e. for every \(\omega\) as Lebesgue Stieltjes measures, with density process \(\lambda\) we have:
\begin{align*}
&\lVert H\rVert_{L^2(M,A)}^2=\mathbb{E}[\int_0^\infty (H_s)^2\diff [M]_s]+\mathbb{E}[\left(\int_0^\infty |H_s||\lambda_s|\diff[M]_s\right)^2]\\&\leq\mathbb{E}[\int_0^\infty (H_s)^2\diff [M]_s+\mathbb{E}[\int_0^{\infty}(H_s)^2\diff [M]_s\cdot\int_0^{\infty}(\lambda_s)^2\diff [M]_s]
\\&
\leq\mathbb{E}[\int_0^\infty (H_s)^2\diff [M]_s+\mathbb{E}[\left(\int_0^{\infty}(H_s)^2\diff [M]_s\right)^{p}]^{1/p}\cdot\mathbb{E}[\left(\int_0^{\infty}(\lambda_s)^2\diff [M]_s\right)^{q}]^{1/q}
\;,
\end{align*}
where we used the fact that if \(A<<[M]\) with density \(\lambda\) then \(|A|<<[M]\) with density \(|\lambda|\). Furthermore we used the Kunita-Watanabe inequality in the second step and Hölder's inequality in the third step. In the case \(p=1,q=\infty\) the above reads
\begin{equation}
    \lVert H\rVert_{L^2(M,A)}\leq\lVert H\rVert_{L^2(M)}\left(1+\operatorname*{ess\sup}\left(\int_0^{\infty}(\lambda_s)^2\diff [M]_s\right)\right)\;.
\end{equation}
Notice \(\int_0^{\infty}(\lambda_s)^2\diff [M]_s=\int_0^{\infty}|\lambda_s|\diff A_s\leq \int_0^{\infty}|\lambda_s|\diff |A|_s\). We impose the additional assumption
\begin{equation}
    \operatorname*{ess\sup}\int_0^{\infty}|\lambda_s|\diff |A|_s<\infty\;.
\end{equation}
Then, because we have the trivial inequality
\begin{align}\label{eq:essential boundedness assumption}
    \mathbb{E}[\int_0^{\infty}H_s^2\diff[M]_s]\leq  \mathbb{E}[\int_0^{\infty}H_s^2\diff[M]_s]+  \mathbb{E}[\left(\int_0^{\infty}|H_s||\lambda_s|\diff[M]_s\right)^2]=\lVert H\rVert_{L^2(M,A)}
\end{align}
we have that the norms on \(L^2(M)\) and \(L^2(M,A)\),under the condition that \(X = X_0+M+A\) with \(A<<[M]\) and \(M\in\mathcal{H}^{p+\epsilon}_0\) for some \(p\in[0,\infty)\) and \(\epsilon>0\), along with the additional assumption \cref{eq:essential boundedness assumption}, are equivalent. Because both spaces \(L^2(M)\) and \(L^2(M,A)\) are subspaces of the space of \(\mathcal
P\)-measurable functions in, where \(\mathcal{P}\) is the predictable \(\sigma\)-algebra on \(\Omega\times[0,\infty)\), they are defined through their norms; the spaces are equal:
\begin{equation}
    L^2(M) =L^2(M,A)\;.
\end{equation}
Hence, density of \(\mathcal{E}(\psi)\) in \(L^2(M)\) is obviously equivalent to density of \(\mathcal{E}(\psi)\) in \(L^2(M,A)\) under the stated assumptions on \(M\) and \(X\). Hence, for every \(H\in L^2(M,A)\) there exists \((H_n)_{n\in\mathbb{N}}\subset \mathcal{E}(\psi)\) such that \(\lim_{n\to\infty}\lVert H-H_n\rVert_{L^2(M,A)}=0\). Then we have, by the observation \cref{eq:bounding the L^2 norm of the stochastic integral}, with \(T=\infty\) and \(H_s\coloneqq (H-H_n)_s\)
\begin{equation}
    \lVert (H-H^n)\bullet X\rVert_{L^2}
    \lesssim\lVert H-H^n\rVert_{L^2(M,A)}\;,
\end{equation}
thus
\begin{equation}
    \lim_{n\to\infty}\lVert (H-H^n)\bullet X\rVert_{L^2}=0
\end{equation}
\bigbreak
The two approaches differ in their additional assumption
\begin{enumerate}
    \item \(\operatorname*{ess\sup}(\langle M\rangle_{\infty}(1+\lambda^2)^{*}_{\infty})<\infty\)
    \item  \(\operatorname*{ess\sup}\int_0^{\infty}|\lambda_s|\diff |A|_s<\infty\)
\end{enumerate}
Notice that
\begin{equation}
    \int_0^\infty(\lambda_s)^2\diff [M]_s\leq\int_0^{\infty}(1+\lambda_s^2)\diff[M]_s\leq(1+\lambda_s^2)_{\infty}^*[M]_{\infty}\,
\end{equation}
where we used the non-negativity and monotonicity of the process \([M]\). Hence, (2.) is a weaker assumption than (1.).
\textcolor{blue}{potential problem with this approach. The norms are equivalent but not really norms. First one needs to check that \(\lVert \cdot\rVert_{L^2(M,A)}\) defines at least a subadditiv, homogenous map and then that the process which are identified by \(\lVert X-Y\rVert_{L^2(M,A)}=0\) are also identified in \(L^2(M)\), that is \(\lVert X-Y\rVert_{L^2(M)}=0\). So that the norms are norms on the same quotient space under some equivalence relation identifying processes}
First, notice that \(\lVert H\rVert_{L^2(M,A)}=0\) is subadditive
Under the assumption that \(A<<[M]\) has density \(\lambda\), we have that \(\lVert X-Y\rVert_{L^2(M)}=0\implies\lVert X-Y\rVert_{L^2(M,A)}=0\). This follows from the fact that \(c\lVert X\rVert_{L^2(M)}\leq\lVert X\rVert_{L^2(M,A)}\leq C\lVert X\rVert_{L^2(M)}\) for \(c=1,C=(1+\operatorname*{ess\sup}\int_0^{\infty}|\lambda_s|\diff |A|_s)\). By assumption, \(c,C\) are finite. This implies that under the equivalence relation \(X\sim Y\) for \(X\sim Y\) for \(\lVert X-Y\rVert_{L^2(M,A)}=0\). Hence, they are norms on the same space \(L^2(\Omega\times [0,T],\mathcal{P})/\sim\), and in particular, it makes sense to say that they are equivalent. 
\bigbreak
This proves the following result:
\begin{theorem}[Semimartingale Deep Hedging]\label{thm:general deep hedging}
    Let \(X=X_0+M+A\) be a semimartingale with \(M\in \mathcal{H}_{0}^{2+\epsilon}\) for some \(\epsilon>0\), \(A<<[M]\) with density \(\lambda=(\lambda_s)_{s\in\mathbb{R}_+}\) and with \(\operatorname*{ess\sup}\int_0^{\infty}|\lambda_s|\diff |A|_s<\infty\) \. Let \(H\in L^2(\mathbb{P})\) and \((V_0^*,\vartheta^*)\) be the solution to \cref{thm:Hilbert Projection for Mean Sqaure Hedging}. Then there exists \((\vartheta_n)_{n\in\mathbb{N}}\subset\mathcal{E}(\mathcal{\psi})\) such that:
    \begin{equation}
        \min_{\vartheta\in \Theta}\lVert H-V_0^*-\int_{0}^T\vartheta_s\diff X_s\rVert_{L^2}=\inf_{\vartheta\in\mathcal{E}(\psi)}\lVert H-V_0^*-\int_{0}^T\vartheta_s\diff X_s\rVert_{L^2}\;,
    \end{equation}
\end{theorem}
where \(\lVert H-V_0^*-\int_{0}^T\vartheta^*_s\diff X_s\rVert_{L^2}=\min_{\vartheta\in \Theta}\lVert H-V_0^*-\int_{0}^T\vartheta_s\diff X_s\rVert_{L^2}\).
\begin{proof}
    The proof is given above and relies on the fact that in this case \(L^2(M)=L^2(M,A)\). By \parencite[Thm.~5.32]{Arandjelovic2025} \(\mathcal{E}(\psi)\) is dense in \(L^2(M)\) and hence also trivially in \(L^2(M,A)\) as well. The only difference is that the optimal \(V_0^*,H^*\) are for the more general problem with \(X= X_0+M+A\). The proof is the same as in \parencite[THm.~5.35]{Arandjelovic2025} as for any \(\delta>0\) there exists \(U\in\mathcal{E}(\psi)\) such that \(\lVert(U^*\bullet X)_{\infty}-(U\bullet X)_{\infty}\rVert_{L^2}\leq \lVert U^*-U\rVert_{L^2(M,A)}<\delta\) by the calculation in \cref{eq:bounding the L^2 norm of the stochastic integral}.
\end{proof}
One should give examples, if they exist, of finite variation processes \(A\) satisfying the assmuptions:
\begin{itemize}
    \item \(A<<[M]\) with density \(\lambda\)
    \item \(\operatorname*{ess\sup}\int_0^{\infty}|\lambda_s|\diff |A|_s<\infty\)
\end{itemize}

Furthermore we will think of a more general approach generalizing \parencite[Thm.~5.35]{Arandjelovic2025} again with \(X\) a semimartingale and approximants from \(\mathcal{S}(\psi)\). For this we might think again about strategy (1) and if there are sensible conditions on \(\vartheta\) that transfer to conditions on \(\vartheta_n\) such that the respective family of stochastic integrals (more precisely the square processes) become a uniformly integrable family.

\textcolor{blue}{Possible, non-trivial process satisfying assumptions: Let}
\begin{equation}
    S_t = \lambda B_t + \mu t
\end{equation}
Then \(\langle S\rangle_t= \lambda^2 t\) and \(A_t=\mu t\).
\end{document}
% so it automatically matches the exact same layout/preamble.

\section{Proof attempts}

% Put your drafts here. Example skeleton:
%
%\begin{proof}[Proof attempt (Lemma~\ref{lem:...})]
%    ...
%\end{proof}
\section{Extension to semi-martingale case}
\cref{thm:cont-deep-ms-hedging-martingale} was in the realm of classical martingale stochastic calculus, with integrators \(M\in\mathcal{H}^{2}\)(even \(\mathcal{H}^{2+\epsilon}\) and integrands \(\vartheta\in L^2(M)\). In light of \cref{thm:Approx of stoch. Integrals}, we aim to prove a similar result as in \cref{thm:cont-deep-ms-hedging-martingale} for the more general case where \(M\) is a semimartingale. Following \parencite{Schweizer2001GuidedTourQuadraticHedging} we restrict the class of possible strategies to \(\vartheta_{GLP}\) such that \(G_T(\vartheta_{GLP})\) is a closed subspace in \(L^2(\mathbb{P})\). In this case, we know that the minimization problem has a solution, as \(\mathcal{A}=\mathbb{R}+G_T(\Theta_{GLP})\) is a closed subspace in \(L^2(\mathbb{P})\).; this means there exists an optimal trading strategy \((V_0^*,\vartheta^*)\in\mathcal{A}=\mathbb{R}+G_T(\Theta_{GLP})\) (\textcolor{blue}{I guess we don't even need strong conditions to get \(\mathcal{A}\) closed as we don't need explicit expression of the optimal trading strategy at first, hence we can directly project on \(\bar{\mathcal{A}}\) which is closed in \(L^2(\mathbb{P})\). In that line of thought, any \(\Theta\subset\Theta_2\) is sufficient for the existence of a solution in \(\bar{\mathcal{A}}\)}.) Then we have the following consequence of the Hilbert Projection theorem
\begin{theorem}[Hilbert Projection for Mean Sqaure Hedging]\label{thm:Hilbert Projection for Mean Sqaure Hedging}
    Let \(\Theta\subset\Theta_2\) and \(\mathcal{A}=\mathbb{R}+G_T(\Theta)\subset L^2(\mathbb{P})\). Then the \(L^2\)-closure \(\bar{\mathcal{A}}\) is a closed subspace of \(L^2(\mathbb{P})\). Hence, there exists \((V_0^*,\vartheta^*)\) such that 
    \begin{equation}
   \lVert H-V_0^*-\int_{0}^T\vartheta^*_s\diff X_s\rVert_{L^2}= \inf_{(c,\vartheta)\in\bar{\mathcal{A}}}\lVert H-c-\int_{0}^T\vartheta_s\diff X_s\rVert_{L^2}
    \end{equation}
    and the solution is unique. 
\end{theorem}
\begin{remark}
    Notice, we can study the related problem where we fix \(V_0\) (corresponds to the situation where the price \(V_0\) of the liability \(H\) might be determined by the market) and study
    \begin{equation}
        \inf_{\vartheta\in\Theta}\lVert H-V_0-\int_{0}^T\vartheta_s\diff X_s\rVert_{L^2}
    \end{equation}
    Again we can consider the the \(L^2\) closure of \(G_T(\Theta)=\left\{\int_0^T \vartheta_s\diff X_s:\vartheta\in\Theta\right\}\) and use the Hilbert Projection theorem on this closed subspace of \(L^2\).
\end{remark}
Then, crucially, notice that from \cref{thm:Approx of stoch. Integrals} we obtain the convergence in probability of the final values of the respective integral processes. Indeed, the \(\mathbb{P}\)-\(\mathbb{P}\)-ucp convergence of the integral processes, restricted to a finite time horizon \(T\), reads as:
\begin{equation}
    \sup_{0\leq t\leq T}\left|\int_{0}^t\vartheta_s^n\diff X_s-\int_{0}^t\vartheta^*_s\diff X_s\right|\to 0\quad\text{in Probability}
\end{equation}
where \(\{\vartheta^n\}_{n\in\mathbb{N}}\in\mathcal{S}(\psi)\) is a minimizing sequence from \cref{thm:Approx of stoch. Integrals}
\\\\
In particular, this implies 
\begin{equation}
    \int_{0}^T\vartheta_s^n\diff X_s\xrightarrow{\mathbb{P}}\int_{0}^T\vartheta^*_s\diff X_s\;,
\end{equation}
where convergence holds in probability. Imposing the additional assumption that \(\left\{\left|\int_{0}^T\vartheta_s^n\diff X_s\right|^2\right\}_{n\in\mathbb{N}}\) is a uniformly integrable family of random variables  leads to the \(L^2\) convergence\footnote{In particular, the condition also implies that the random variable \(\int_{0}^T\vartheta_s^n\diff X_s\in L^2(\mathbb{P)}\quad\forall n\in\mathbb{N}\)}: 
\begin{equation}
\int_{0}^T\vartheta_s^n\diff X_s \xrightarrow{L^2(\mathbb{P})} \int_{0}^T\vartheta^*_s\diff X_s\;.
\end{equation}
We do NOT want the uniform integrability of \(\left\{\left|\int_{0}^T\vartheta_s^n\diff X_s\right|^2\right\}_{n\in\mathbb{N}}\) to follow from conditions on a sequence of trading strategies. Indeed, \cref{thm:Approx of stoch. Integrals} asserts the existence of a sequence \(\{\vartheta^n\}_{n\in\mathbb{N}}\subset\mathcal{S}(\psi)\). We don't want a theorem that says something like: 
\\\\
"If that sequence happens to be square uniformly integrable, then it follows...."
\\\\
Hence, if necessary, we shall impose conditions only on the integrator \(X\). Recall that for \(\vartheta_n\in\mathcal{S}(\psi)\) that is \(\vartheta^n_t=\varphi_0\mathbbm{1}_{\{0\}}+\sum_{i=1}^l\varphi_i^n\mathbbm{1}_{\tau_i,\tau_{i+1}}\) and \(\varphi_i^{(n)}=f^{(n)}(Y^{\tau_i}_{t_1},\dots,Y^{\tau_i}_{t_r})\) where \(f\in\mathcal{NN}((\mathbb{R}^{d\times r})^*,\psi)\)\footnote{the parameter \(r\in\mathbb{N}\) is not further specified. It relies to the existence result of \parencite[Thm.~5.23.]{Arandjelovic2025}}, the stochastic integral takes the form \(\int_0^T\vartheta^n_s\diff X_s = \sum_{i=1}^k \varphi^{(n)}_i(X_{\tau_i\wedge T}-X_{\tau_{i+1}\wedge T})\). Note that \(\left|\int_0^T\vartheta^n_s\diff X_s\right|^2 \leq \sum_{i=1}^k \left|\varphi^{(n)}_i\right|^2\left|(X_{\tau_i\wedge T}-X_{\tau_{i+1}\wedge T})\right|^2\). But there can't exist a uniform bound (in n) on \(|\varphi_i^n|\) as this would imply that neural networks can't approximate constants above this bound which is absurd. Also, notice for \(\tau<\sigma\) two stopping times, that \(f_n = c_n\mathbbm{1}_{[\tau,\sigma)}\in\mathcal{S}(\psi)\) for all \(c_n\in\mathbb{R}\). Take \(c_n\to\infty\) as \(n\to\infty\) a real deterministic sequence, then \(\int_0^T(f_n)_s\diff X_s=c_n(X_{\tau\wedge T}-X_{\sigma\wedge T})\), and hence \(\mathbb{E}[\left|\int_0^T(f_n)_s\diff X_s\right|^2]=\mathbb{E}[\left|c_n\right|^2\left|(X_{\tau\wedge T}-X_{\sigma\wedge T})\right|^2]\). If this is to be UI, then in particular, it is bounded in \(L^1\). But \(\sup_{n\in\mathbb{N}}\mathbb{E}[c_n^2(X_{\tau\wedge T}-X_{\sigma\wedge T})^2]=(\sup_{n\in\mathbb{N}}c_n^2)\mathbb{E}[(X_{\tau\wedge T}-X_{\sigma\wedge T})^2]<\infty\) by assumption. But \(\sup_{n\in\mathbb{N}}c_n^2=\infty\) hence we must have \(\mathbb{E}[(X_{\tau\wedge T}-X_{\sigma\wedge T})^2]=0\) which implies \(X_{\tau\wedge T}-X_{\sigma\wedge T}=0\quad\mathbb{P}\)-a.s. As this holds for all \(\mathbb{F}\)-stopping times \(\tau,\sigma\), in particular deterministic ones, we have \(X\equiv 0\) almost surely. Hence, there can't be a (non-trivial) condition on \(X\) making the family \(\left\{\left|\int_{0}^T\vartheta_s^n\diff X_s\right|^2\right\}_{n\in\mathbb{N}}\) UI, as the latter being UI for any sequence \((\vartheta^n)_{n\in\mathbb{N}}\) implies that \(X\equiv 0\).
\bigbreak
\textcolor{blue}{As an initial observation, recall that \(\int_{0}^T\vartheta_s^n\diff X_s\xrightarrow{\mathbb{P}}\int_{0}^T\vartheta^*_s\diff X_s\;\) implies that there exists a subsequence \(\left\{\int_{0}^T\vartheta_s^{n(k)}\diff X_s\right\}_{k\in\mathbb{N}}\) that converges \(\mathbb{P}\)-a.s. to \(\int_{0}^T\vartheta_s^{*}\diff X_s\). Hence, on a full measure set, it is (uniformly) bounded as a real converging sequence (assuming that \(\int_{0}^T\vartheta_s^{*}\diff X_s<\infty\) on that full measure set). If it is only \(\mathbb{P}\)-a.s. finite, take the intersection of the respective full measure sets to reach the same conclusion.) Hence, there exists a positive constant \(M\in\mathbb{R}_+\) such that \(\left|\int_{0}^T\vartheta_s^{n(k)}\diff X_s\right|^2<M^2\quad \mathbb{P}\)-a.s. for all \(k\in\mathbb{N}\). But \(M^2\) is integrable; hence, the family \(\left\{\left|\int_{0}^T\vartheta_s^{n(k)}\diff X_s\right|^2\right\}_{k\in\mathbb{N}}\) is uniformly integrable.\footnote{If \(H\geq |X_n|\) for all \(n\in\mathbb{N}\), we have \(\mathbb{E}[|X|\mathbbm{1}_{|X|\geq K}]\leq\mathbb{E}[H\mathbbm{1}_{H\geq K}]\) for any \(K>0\). As the left hand side going to 0 as \(K\to\infty\) it gives us the uniform integrability of \(\{X_n\}_{n\in\mathbb{N}}\)}. Thus, we get }
\\\\
\textcolor{red}{This conclusion is WRONG. Notice that the bound \(M = M(\omega)\) and hence doesn't define a random variable \(H\in L^2\)(at least not because it is constant almost surely). Indeed the argument used, goes as follows: For every \(\omega\) the sequence \(\left(\int_{0}^T\vartheta_s^{n(k)}\diff X_s\right)(\omega)\) is bounded as a real converging sequence, hence \(\left|\left(\int_{0}^T\vartheta_s^{n(k)}\diff X_s\right)(\omega)\right|\leq M(\omega)\).}
\\\\
\begin{equation}
\int_{0}^T\vartheta_s^{n(k)}\diff X_s \xrightarrow{L^2(\mathbb{P})} \int_{0}^T\vartheta^*_s\diff X_s\;\text{ as } k\to\infty
\end{equation}
Hence, restricting ourselves to this subsequence (which obviously still converges in probability to the same limit), we have 
\begin{align*}
    &\inf_{k\in\mathbb{N}}\Vert[H-V_0^*-\int_{0}^T\vartheta_s^{n(k)}\diff X_s]\rVert_{L^2}\leq \liminf_{k\to\infty}\lVert H-V_0^*-\int_{0}^T\vartheta_s^{n(k)}\diff X_s]\rVert_{L^2}\\&\leq \liminf_{k\to\infty}\lVert H-V_0^*-\int_{0}^T\vartheta_s^{*}\diff X_s\rVert_{L^2}+\lVert \int_{0}^T\vartheta_s^{*}\diff X_s -\int_{0}^T\vartheta_s^{n(k)}\diff X_s\rVert_{L^2}
    \\&=\lVert H-V_0^*-\int_{0}^T\vartheta_s^{*}\diff X_s\rVert_{L^2}
\end{align*}
where \(\lVert H-V_0^*-\int_{0}^T\vartheta_s^{*}\diff X_s\rVert_{L^2}=\inf_{\vartheta\in\Theta}\lVert H-V_0^*-\int_{0}^T\vartheta_s\diff X_s\rVert_{L^2}\).
By \cref{rem:inf and sequences}, the existence of this sequence is equivalent to a more general statement:
\begin{theorem}[Semimartingale Deep Hedging]\label{thm:general deep hedging}
    Let \(X\) be a semimartingale and \(\Theta\subset\Theta_2\). Let \((V_0^*,\vartheta^*)\) be the solution to \cref{thm:Hilbert Projection for Mean Sqaure Hedging}. Then there exists \((\vartheta_n)_{n\in\mathbb{N}}\subset\mathcal{S}(\mathcal{\psi})\) such that:
    \begin{equation}
        \min_{\vartheta\in \Theta}\lVert H-V_0^*-\int_{0}^T\vartheta_s\diff X_s\rVert_{L^2}=\inf_{\vartheta\in\mathcal{S}(\psi)}\lVert H-V_0^*-\int_{0}^T\vartheta_s\diff X_s\rVert_{L^2}\;,
    \end{equation}
\end{theorem}
where \(\lVert H-V_0^*-\int_{0}^T\vartheta^*_s\diff X_s\rVert_{L^2}=\min_{\vartheta\in \Theta}\lVert H-V_0^*-\int_{0}^T\vartheta_s\diff X_s\rVert_{L^2}\).
\\\\
\newpage
\underline{Possible strategies}:
In order to get deep-hedging result similar to \cref{thm:general deep hedging} we need 
\begin{equation}
    \int_0^T\vartheta^n_s\diff X_s \overset{L^2}{\to} \int_0^T\vartheta_s\diff X_s
\end{equation}
For an approximating sequence \((\vartheta_n)_{n\in\mathbb{N}}\subset\mathcal{S}(\psi)\) of \(\vartheta\in\tilde{\Theta}\), where the sense of approximation is to be determinded and the set \(\tilde{\Theta}\) such that \(\mathcal{S}(\psi)\subset\tilde{\Theta}\). Assuming this is done, we have
\begin{align*}
  &\inf_{\vartheta\in\mathcal{S}(\psi)}\lVert H-V_0^*-\int_0^T \vartheta_s^n\diff X_s\rVert_{L^2}\\&\leq
    \inf_{n\in\mathbb{N}}\lVert H-V_0^*-\int_0^T \vartheta_s^n\diff X_s\rVert_{L^2}\leq \liminf_{n\to\infty}\lVert H-V_0^*-\int_0^T \vartheta_s^n\diff X_s\rVert_{L^2}
    \\&\leq \liminf_{n\to\infty}\lVert H-V_0^*-\int_0^T \vartheta_s^*\diff X_s\rVert_{L^2}+\lVert \int_0^T \vartheta_s^n\diff X_s-\int_0^T \vartheta_s^*\diff X_s\rVert_{L^2}\\&=\lVert H-V_0^*-\int_0^T \vartheta_s^*\diff X_s\rVert_{L^2}\;,
\end{align*}
where \(\lVert H-V_0^*-\int_0^T \vartheta_s^*\diff X_s\rVert_{L^2}=\argmin_{\vartheta\in\Theta}\lVert H-V_0^*-\int_0^T \vartheta_s^*\diff X_s\rVert_{L^2}\) for some set \(\Theta\subset\Theta_2\). On the other hand we have \(\inf_{\vartheta\in\tilde{\Theta}}\lVert H-V_0^*-\int_0^T \vartheta_s^n\diff X_s\rVert_{L^2}\leq\inf_{\vartheta\in\mathcal{S}(\psi)}\lVert H-V_0^*-\int_0^T \vartheta_s^n\diff X_s\rVert_{L^2}\), because \(\mathcal{S}(\psi)\subset\tilde{\Theta}\). Thus we get equality. \\
We consider two possible strategies.
\bigbreak
(1): From \parencite[Thm~5.29.]{Arandjelovic2025} we have \(\mathcal{S}(\psi)\) is dense in \(\mathcal{S}\) which is dense in \(\mathcal{L}\), the space of càglàd processes (see Protter for this). Then for \(\vartheta_n\overset{\mathbb{P}-ucp}{\to}\vartheta\) where \(\vartheta\in \mathcal{L}\), we have 
\begin{align*}
    \int \vartheta_s^n\diff X_s\overset{\mathbb{P}-ucp}{\to}\int \vartheta_s\diff X_s\implies\int_0^T \vartheta_s^n\diff X_s\overset{\mathbb{P}}{\to}\int_0^T\vartheta_s\diff X_s
\end{align*}
and 
\begin{align*}
    &\left(\vartheta_s^n\diff X_s\overset{\mathbb{P}}{\to}\int_0^T\vartheta_s\diff X_s \right) + \left(\left\{\left|\int_{0}^T\vartheta_s^{n}\diff X_s\right|^2\right\}_{n\in\mathbb{N}}\text{uniformly integrable}\right)\\&\implies \left(\int_0^T\vartheta_s^n\diff X_s\overset{L^2}{\to}\int_0^T\vartheta_s\diff X_s \right)
\end{align*}
WARNING: Notice that for sequences of the form \(\vartheta_n=c_n\cdot\mathbbm{1}_{[\tau,\sigma)}\in\mathcal{S}(\psi)\) for \(\lim_{n\to\infty}c_n=\infty\), requiring that \(\left\{\left|\int_{0}^T\vartheta_s^{n}\diff X_s\right|^2\right\}_{n\in\mathbb{N}}\) is uniformly integrable implies \(X=0\) as a process. This is a pathological example as it doesn't converge to an interesting limiting process \(\vartheta\), which we assume. This just shows that  \(\left\{\left|\int_{0}^T\vartheta_s^{n}\diff X_s\right|^2\right\}_{n\in\mathbb{N}}\) can't be UI for all \((\vartheta_n)_{n\in\mathbb{N}}\subset\mathcal{S}(\psi)\). 
\bigbreak
(2): Notice that \(Y_n\overset{L^2-ucp}{\to}Y\) i.e. \(\lim_{n\to\infty}\mathbb{E}[\sup_{0\leq t\leq T}|Y_t^n-Y_t|^2]=0\) implies \(Y^n_T\overset{L^2}{\to}Y_T\). Hence, we want a version of \parencite[Thm~5.29]{Arandjelovic2025} with \(L^2-ucp\) convergence of the integrals. For this, we might want to use the fact that for \(\Phi(\cdot)=(\cdot)^2\), \parencite[Thm~5.28]{Arandjelovic2025} implies that \(\mathcal{S}(\psi)\) is \(L^2-ucp\) dense in \(\mathcal{S}\) and use some kind of \(L^2-ucp\) continuity of the stochastic integral \(H\to H\bullet X\) for well chosen classes for \(H\) and \(X\). An clearly weaker assumption would be to ask for convergence in \(\mathcal{H}^2\) which implies the \(L^2\) convergence of the end-points (Doob). A well known case where we have the continuity of the integral map w.r.t. to the \(L^2\)-convergence of the end-points is the isometry property (Itô isomtrey) of the map 
\begin{align}
    \begin{cases}
        H\to H\bullet M\\
        L^2(M)\to\mathcal{H}^2_0
    \end{cases}
\end{align}
for \(M\in\mathcal{H}^2_{0,loc}\)
Indeed in this case, by \parencite[Thm~5.32]{Arandjelovic2025} we even have that \(\mathcal{E}(\psi)\) is dense in \(L^p(M)\) for \(p\in[1,\infty)\) if \(M\in\mathcal{H}_0^{p+\epsilon}\) for some \(\epsilon>0\). (Then we also have that \(\mathcal{S}(\psi)\) is dense in \(L^p(M)\)). Hence we have 
\begin{equation}
    \lVert H\bullet M\rVert_{\mathcal{H}^2}=\lVert H\rVert_{L^2(M)}
\end{equation}
Using the sequence \(H_n=\vartheta_n-\vartheta\) and the fact that convergence in \(\mathcal{H}^2\) implies convergence in \(L^2-ucp\) (by Doob's maximal inequality, which is applicable because \(H\bullet M\) is a square integrable martingale in this case) we recover \parencite[Thm~5.35]{Arandjelovic2025}.
\\\\
To get a slight extension of \parencite[Thm~5.35]{Arandjelovic2025}, we notice that we don't need the full isometry property; rather, we need the boundedness (continuity) of the map \(H\to H \bullet M\) and the corresponding density of \(\mathcal{S}(\psi)\). Hence, for a semimartingale 
\begin{equation}
    X = X_0 + A + M
\end{equation}
for \(A\) a finite-variation process and \(M\in\mathcal{H}^2_{0,loc}\) (the addition of \(A\) destroys the Itô-isometry property) we try to bound 
\begin{equation}
     \lVert \left(H\bullet X\right)_T\rVert_{L^2}\leq K_X\lVert H\rVert_{L^2(M)}
\end{equation}
for some constant \(K_X = K(X)\). By \(|a+b|^2\leq 2|a|^2+2|b|^2\) it is enough to bound \(\lVert H\bullet A\rVert_{\mathcal{H}^2}\) and \(\lVert H\bullet M\rVert_{\mathcal{H}^2}\) both by \(\lVert H\rVert_{L^2(M)}\) and some multiplicative constant. Notice that for \(M\in\mathcal{H}^2_{0,loc}\) and \( H \in L^2(M)\), \(H\bullet M\) is a square integrable martingale; hence, we can bound \(\lVert H\bullet M\rVert_{\mathcal{H}^2}=\lVert (H\bullet M)_T\rVert_{L^2}\leq 4\lVert H\rVert_{L^2}\) by Doob. For the finite variation part, we might need assumptions on \(A\) (and hence on \(X\)).
\\
Based on this oberservations we consider the following approach. Consider a new space \(L^2(M,A)\) for a semimartingale \(X=X_0+M+A\) where \(A\) is a finite-variation process and \(M\in\mathcal{H}^2_{0,loc}\). Define it by \footnote{Note that in our case \([X]=[M]\) and \(L^2(M,A)\subset L^2(M)\), so \(\int_0^T\vartheta_s\diff M_s\in\mathcal{H}_0^2\); hence, \(\left(\int \vartheta_s\diff M_s\right)^2-[\int \vartheta_s\diff M_s]\) is a martingale starting at 0, and \([\int\vartheta_s\diff M_s]=\int \vartheta_s^2\diff [M]_s\). Hence, \(\mathbb{E}[\left(\int_0^T \vartheta_s\diff M_s\right)^2]=\mathbb{E}[\int_0^T \vartheta_s^2\diff [M]_s]\)}
\begin{align*}
H\in L^2(M,A) \iff \left(\mathbb{E}[\int_0^T H_s^2\diff [X]_s]+\mathbb{E}[\left(\int_0^T |H|_s\diff |A|_s\right)^2]\right)^{\frac{1}{2}}\coloneqq\lVert H\rVert_{L^2(M,A)}<\infty
\end{align*}\footnote{Inspired by the fact that for \(\int H_s\diff X_s = \int H_s\diff M_s+\int H_s\diff A_s\) and \(\mathbb{E}[\left(\int H_s\diff M_s+\int H_s\diff A_s\right)^2]\leq 2\mathbb{E}[\left(\int H_s \diff M_s\right)^2]+2\mathbb{E}[\left(\int H_s \diff A_s\right)^2]\leq2\cdot 4\mathbb{E}[\left(\int (H_s)^2 \diff [M]_s\right)^2]+2\mathbb{E}[\left(\int H_s \diff A_s\right)^2]\) from \(|a+b|^2\leq 2|a|^2+2|b|^2\) and BDG with \(C_2 = 4\).}\textcolor{blue}{One would still need to check that \(\lVert H\rVert_{L^2(M,A)}\) defines a norm}
where in our case \([X_0 + A + M]=[M]\). With this we have 
\begin{equation}
    \lVert\int_0^T\vartheta_s\diff X_s\rVert_{L^2}\leq K_X\lVert \vartheta \rVert_{L^2(M,A)}
\end{equation}
Indeed,
\begin{align*}\label{eq:bounding the L^2 norm of the stochastic integral}
    &\mathbb{E}[\left(\int_0^T H_s\diff X_s\right)^2]=\mathbb{E}[\left(\int_0^T H_s\diff M_s+\int_0^T H_s\diff A_s\right)^2]
    \\&
    \leq 2\mathbb{E}[\left(\int_0^T H_s \diff M_s\right)^2]+2\mathbb{E}[\left(\int_0^T H_s \diff A_s\right)^2]
    \\&
    \leq2\cdot 4\mathbb{E}[\int_0^T (H_s)^2 \diff [M]_s]+2\mathbb{E}[\left(\int_0^T H_s \diff A_s\right)^2]\leq
    \\&
     \leq2\cdot 4\mathbb{E}[\int_0^T (H_s)^2 \diff [M]_s]+2\mathbb{E}[\left(\int_0^T |H|_s \diff |A|_s\right)^2]\leq
    \lVert H\rVert_{L^2(M,A)}
\end{align*}
from \(|a+b|^2\leq 2|a|^2+2|b|^2\) and BDG with \(C_2 = 4\)\footnote{\(\mathbb{E}[\left(\int_0^TH_s\diff M_s\right)^2]\leq\mathbb{E}[\sup_{0\leq t\leq T}\left(\int_0^t H_s\diff M_s\right)^2]\leq C_2\mathbb{E}[\int_0^T(H_s)^2\diff [M]_s]\)}. Then we have
\begin{align}
    &\vartheta_n\overset{L^2(M,A)}{\to}\vartheta\implies\int_0^T\vartheta_s^n\diff X_s \overset{L^2}{\to}\int_0^T\vartheta_s\diff X_s
\end{align}
In this approach, we need to show that \(\lVert\cdot \rVert_{L^2(M,A)}\) is indeed a norm (or a semi-norm), that \(\mathcal{S}(\psi)\) is dense in \(L^2(M,A)\), and whether \(L^2(M,A)\) is rich enough to be interesting.
\section{Proof Attempt}
We present the first version of a generalization of \parencite[Thm~5.35]{Arandjelovic2025}. From \parencite[Thm.~5.32.]{Arandjelovic2025} we know that \(\mathcal{E}(\psi)\) is dense in \(L^p(M)\) for \(p\in[1,\infty)\) and \(M\in\mathcal{H}_0^{p+\epsilon}\) for some \(\epsilon>0\). 
Now consider a semimartingale \(X=X_0+M+A\), where \(M\in\mathcal{H}_0^{p+\epsilon}\) and \(A\) are adapted processes with RCLL trajectories, \(A_.(\omega)\) of finite variation and \(A_0=0\). Additionally, assume the following:
\begin{itemize}
    \item \(A<<\langle M\rangle\) path-wise (i.e. for every \(\omega\) as Lebesgue Stieltjes measures, with density process \(\lambda\).
    \item \(\operatorname*{ess\sup}\left(\langle M\rangle_{\infty}(\lambda^2+1)^*_{\infty}\right)<\infty\) where \((\lambda^2+1)^*_{\infty}=\sup_{t\geq 0}(\lambda^2+1)\).\footnote{Replace \(\infty\) by \(T\) in the finite horizon setting.}
\end{itemize}
Then we have, by density, of \(\mathcal{E}(\psi)\) in \(L^2(M)\) that
\begin{equation}
    \forall H\in L^2(M), \exists (H_n)_{n\in\mathbb{N}}\subset\mathcal{E}(\psi):\lim_{n\to\infty}\lVert H-H_n\rVert_{L^2(M)}=0\;.
\end{equation}
We can bound the \(L^2\) difference of the integrals by 
\begin{align}
    &\mathbb{E}[(H\bullet X)_{\infty}^2]\leq 2\left((\mathbb{E}[(H\bullet \langle M\rangle)_{\infty}^2]+\mathbb{E}[(H\bullet A)_{\infty}^2]\right)\\
    &\leq 2\left((\mathbb{E}[\langle M\rangle_{\infty}(H^2\bullet \langle M\rangle)_{\infty}]+\mathbb{E}[(H\lambda\bullet \langle M\rangle)_{\infty}^2]\right)\\
    &\leq2\left((\mathbb{E}[\langle M\rangle_{\infty}(H^2\bullet \langle M\rangle)_{\infty}]+\mathbb{E}[\langle M\rangle_{\infty}(H^2\lambda^2\bullet \langle M\rangle)_{\infty}]\right)\\
    &=2\mathbb{E}[\langle M\rangle_{\infty}(H^2(1+\lambda^2)\bullet\langle M\rangle)_{\infty}]\\
    &\leq 2\mathbb{E}[\langle M\rangle_{\infty}(1+\lambda^2)^{*}_{\infty}(H^2\bullet\langle M\rangle)_{\infty}]\\
    &\leq2\operatorname*{ess\sup}(\langle M\rangle_{\infty}(1+\lambda^2)^{*}_{\infty})\cdot\mathbb{E}[H^2\bullet \langle M\rangle]\\
    &=C(X)\lVert H\rVert_{L^2(M)}^2\;.
\end{align}
We used Cauchy Schwarz in step 2 and 3, along with the absolute continuity of \(A<<\langle M\rangle\). In step 5, we used the Hölder inequality with \(p=1,q=\infty\).\\
By assumption \(C(X)<\infty\), hence substituting for \(H\), \(Z^n=H-H^n\), we get that 
\begin{equation}
    \lim_{n\to\infty}\mathbb{E}[(Z^n\bullet X)_{\infty}^2]\leq\lim_{n\to\infty}C(X)\lVert H\rVert_{L^2(M)}^2=0
\end{equation}
By the argument given earlier, this allows for a deep mean-variance hedging result for a subclass of semimartingales.
\section{Different Approach}
We present a slightly different approach based on the observation \cref{eq:bounding the L^2 norm of the stochastic integral}. Notice that under the assumption that \(A<<\langle M\rangle\) path-wise (i.e. for every \(\omega\) as Lebesgue Stieltjes measures, with density process \(\lambda\) we have:
\begin{align*}
&\lVert H\rVert_{L^2(M,A)}^2=\mathbb{E}[\int_0^\infty (H_s)^2\diff [M]_s]+\mathbb{E}[\left(\int_0^\infty |H_s||\lambda_s|\diff[M]_s\right)^2]\\&\leq\mathbb{E}[\int_0^\infty (H_s)^2\diff [M]_s+\mathbb{E}[\int_0^{\infty}(H_s)^2\diff [M]_s\cdot\int_0^{\infty}(\lambda_s)^2\diff [M]_s]
\\&
\leq\mathbb{E}[\int_0^\infty (H_s)^2\diff [M]_s+\mathbb{E}[\left(\int_0^{\infty}(H_s)^2\diff [M]_s\right)^{p}]^{1/p}\cdot\mathbb{E}[\left(\int_0^{\infty}(\lambda_s)^2\diff [M]_s\right)^{q}]^{1/q}
\;,
\end{align*}
where we used the fact that if \(A<<[M]\) with density \(\lambda\) then \(|A|<<[M]\) with density \(|\lambda|\). Furthermore we used the Kunita-Watanabe inequality in the second step and Hölder's inequality in the third step. In the case \(p=1,q=\infty\) the above reads
\begin{equation}
    \lVert H\rVert_{L^2(M,A)}\leq\lVert H\rVert_{L^2(M)}\left(1+\operatorname*{ess\sup}\left(\int_0^{\infty}(\lambda_s)^2\diff [M]_s\right)\right)\;.
\end{equation}
Notice \(\int_0^{\infty}(\lambda_s)^2\diff [M]_s=\int_0^{\infty}|\lambda_s|\diff A_s\leq \int_0^{\infty}|\lambda_s|\diff |A|_s\). We impose the additional assumption
\begin{equation}
    \operatorname*{ess\sup}\int_0^{\infty}|\lambda_s|\diff |A|_s<\infty\;.
\end{equation}
Then, because we have the trivial inequality
\begin{align}\label{eq:essential boundedness assumption}
    \mathbb{E}[\int_0^{\infty}H_s^2\diff[M]_s]\leq  \mathbb{E}[\int_0^{\infty}H_s^2\diff[M]_s]+  \mathbb{E}[\left(\int_0^{\infty}|H_s||\lambda_s|\diff[M]_s\right)^2]=\lVert H\rVert_{L^2(M,A)}
\end{align}
we have that the norms on \(L^2(M)\) and \(L^2(M,A)\),under the condition that \(X = X_0+M+A\) with \(A<<[M]\) and \(M\in\mathcal{H}^{p+\epsilon}_0\) for some \(p\in[0,\infty)\) and \(\epsilon>0\), along with the additional assumption \cref{eq:essential boundedness assumption}, are equivalent. Because both spaces \(L^2(M)\) and \(L^2(M,A)\) are subspaces of the space of \(\mathcal
P\)-measurable functions in, where \(\mathcal{P}\) is the predictable \(\sigma\)-algebra on \(\Omega\times[0,\infty)\), they are defined through their norms; the spaces are equal:
\begin{equation}
    L^2(M) =L^2(M,A)\;.
\end{equation}
Hence, density of \(\mathcal{E}(\psi)\) in \(L^2(M)\) is obviously equivalent to density of \(\mathcal{E}(\psi)\) in \(L^2(M,A)\) under the stated assumptions on \(M\) and \(X\). Hence, for every \(H\in L^2(M,A)\) there exists \((H_n)_{n\in\mathbb{N}}\subset \mathcal{E}(\psi)\) such that \(\lim_{n\to\infty}\lVert H-H_n\rVert_{L^2(M,A)}=0\). Then we have, by the observation \cref{eq:bounding the L^2 norm of the stochastic integral}, with \(T=\infty\) and \(H_s\coloneqq (H-H_n)_s\)
\begin{equation}
    \lVert (H-H^n)\bullet X\rVert_{L^2}
    \lesssim\lVert H-H^n\rVert_{L^2(M,A)}\;,
\end{equation}
thus
\begin{equation}
    \lim_{n\to\infty}\lVert (H-H^n)\bullet X\rVert_{L^2}=0
\end{equation}
\bigbreak
The two approaches differ in their additional assumption
\begin{enumerate}
    \item \(\operatorname*{ess\sup}(\langle M\rangle_{\infty}(1+\lambda^2)^{*}_{\infty})<\infty\)
    \item  \(\operatorname*{ess\sup}\int_0^{\infty}|\lambda_s|\diff |A|_s<\infty\)
\end{enumerate}
Notice that
\begin{equation}
    \int_0^\infty(\lambda_s)^2\diff [M]_s\leq\int_0^{\infty}(1+\lambda_s^2)\diff[M]_s\leq(1+\lambda_s^2)_{\infty}^*[M]_{\infty}\,
\end{equation}
where we used the non-negativity and monotonicity of the process \([M]\). Hence, (2.) is a weaker assumption than (1.).
\textcolor{blue}{potential problem with this approach. The norms are equivalent but not really norms. First one needs to check that \(\lVert \cdot\rVert_{L^2(M,A)}\) defines at least a subadditiv, homogenous map and then that the process which are identified by \(\lVert X-Y\rVert_{L^2(M,A)}=0\) are also identified in \(L^2(M)\), that is \(\lVert X-Y\rVert_{L^2(M)}=0\). So that the norms are norms on the same quotient space under some equivalence relation identifying processes}
First, notice that \(\lVert H\rVert_{L^2(M,A)}=0\) is subadditive
Under the assumption that \(A<<[M]\) has density \(\lambda\), we have that \(\lVert X-Y\rVert_{L^2(M)}=0\implies\lVert X-Y\rVert_{L^2(M,A)}=0\). This follows from the fact that \(c\lVert X\rVert_{L^2(M)}\leq\lVert X\rVert_{L^2(M,A)}\leq C\lVert X\rVert_{L^2(M)}\) for \(c=1,C=(1+\operatorname*{ess\sup}\int_0^{\infty}|\lambda_s|\diff |A|_s)\). By assumption, \(c,C\) are finite. This implies that under the equivalence relation \(X\sim Y\) for \(X\sim Y\) for \(\lVert X-Y\rVert_{L^2(M,A)}=0\). Hence, they are norms on the same space \(L^2(\Omega\times [0,T],\mathcal{P})/\sim\), and in particular, it makes sense to say that they are equivalent. 
\bigbreak
This proves the following result:
\begin{theorem}[Semimartingale Deep Hedging]\label{thm:general deep hedging}
    Let \(X=X_0+M+A\) be a semimartingale with \(M\in \mathcal{H}_{0}^{2+\epsilon}\) for some \(\epsilon>0\), \(A<<[M]\) with density \(\lambda=(\lambda_s)_{s\in\mathbb{R}_+}\) and with \(\operatorname*{ess\sup}\int_0^{\infty}|\lambda_s|\diff |A|_s<\infty\) \. Let \(H\in L^2(\mathbb{P})\) and \((V_0^*,\vartheta^*)\) be the solution to \cref{thm:Hilbert Projection for Mean Sqaure Hedging}. Then there exists \((\vartheta_n)_{n\in\mathbb{N}}\subset\mathcal{E}(\mathcal{\psi})\) such that:
    \begin{equation}
        \min_{\vartheta\in \Theta}\lVert H-V_0^*-\int_{0}^T\vartheta_s\diff X_s\rVert_{L^2}=\inf_{\vartheta\in\mathcal{E}(\psi)}\lVert H-V_0^*-\int_{0}^T\vartheta_s\diff X_s\rVert_{L^2}\;,
    \end{equation}
\end{theorem}
where \(\lVert H-V_0^*-\int_{0}^T\vartheta^*_s\diff X_s\rVert_{L^2}=\min_{\vartheta\in \Theta}\lVert H-V_0^*-\int_{0}^T\vartheta_s\diff X_s\rVert_{L^2}\).
\begin{proof}
    The proof is given above and relies on the fact that in this case \(L^2(M)=L^2(M,A)\). By \parencite[Thm.~5.32]{Arandjelovic2025} \(\mathcal{E}(\psi)\) is dense in \(L^2(M)\) and hence also trivially in \(L^2(M,A)\) as well. The only difference is that the optimal \(V_0^*,H^*\) are for the more general problem with \(X= X_0+M+A\). The proof is the same as in \parencite[THm.~5.35]{Arandjelovic2025} as for any \(\delta>0\) there exists \(U\in\mathcal{E}(\psi)\) such that \(\lVert(U^*\bullet X)_{\infty}-(U\bullet X)_{\infty}\rVert_{L^2}\leq \lVert U^*-U\rVert_{L^2(M,A)}<\delta\) by the calculation in \cref{eq:bounding the L^2 norm of the stochastic integral}.
\end{proof}
One should give examples, if they exist, of finite variation processes \(A\) satisfying the assmuptions:
\begin{itemize}
    \item \(A<<[M]\) with density \(\lambda\)
    \item \(\operatorname*{ess\sup}\int_0^{\infty}|\lambda_s|\diff |A|_s<\infty\)
\end{itemize}

Furthermore we will think of a more general approach generalizing \parencite[Thm.~5.35]{Arandjelovic2025} again with \(X\) a semimartingale and approximants from \(\mathcal{S}(\psi)\). For this we might think again about strategy (1) and if there are sensible conditions on \(\vartheta\) that transfer to conditions on \(\vartheta_n\) such that the respective family of stochastic integrals (more precisely the square processes) become a uniformly integrable family.

\textcolor{blue}{Possible, non-trivial process satisfying assumptions: Let}
\begin{equation}
    S_t = \lambda B_t + \mu t
\end{equation}
Then \(\langle S\rangle_t= \lambda^2 t\) and \(A_t=\mu t\).
\end{document}